% v0.3 (2026-10-02): bounded C7 diagram alignment and current series scope; historical definitions and raw audit codes retained.
% Bridge_Contract_Criterion_v0_de.tex — Gap-Serie Teil 2 (DE)
% Status: v0.2 MAJOR REVISION (2026-08-16) — strikt nach dem Heiler-Bauplan
% der 3-Agenten-Review-Runde 1
% (_results/3agenten_review/Teil2_BridgeContract_HEILER_2026-08-16.md, §§3–9)
% gebaut, wie von _results/3agenten_review/Teil2_BridgeContract_RICHTER_2026-08-16.md
% entschieden.
% Deutsche Fassung von Bridge_Contract_Criterion_v0_en.tex (1:1-Strukturparität).
% v0.2-Änderungen: Instrumentversionierung v1-F / v1-T mit T-/F-Namensraum;
% Component 6 und 7 gehärtet (verbindliche Kompatibilitätsidentitäten;
% adressiertes Ordnungslemma mit H1–H5; expliziter Auditor-Grenze); Ledger
% ersetzt durch den normativen Crosswalk (F-R6a/F-R6b-Aufspaltung,
% F-R4a/F-R4b -> C5, F-R8 als E8 mit Default X); Kalibrierung umbenannt zu einer
% protokoll-eingefrorenen retrospektiven Zwei-Fall-Illustration von v1-T mit der
% vollständigen Dissensmatrix; abc-Befunde neu gewichtet (F-R2a = X, kein
% künstliches C6-Bestehen, mehrere gewichtete Befunde); Joshi = vergleichender
% Diskriminierungsfall, Prasad die einzige akzeptierte Kontrolle; LANA =
% externer Vergleichspunkt; Geltungsbereich reduziert auf EINE
% abc-Instanziierung plus RH/global-only-Ausblick; KI-Offenlegung eingebunden;
% Operator-Blöcke verschoben nach
% _claude-work/T2_REMOVED_COMMENTS_2026-08-16.md.
% v0.1-Herkunft: Strukturfreigabe D-005=A [U 2026-08-15]; P0-Gates G1–G5
% erledigt; sec:related aus t2-s3 (Crosswalk), sec:abc aus t2-s4;
% sec:calibration aus dem protokoll-eingefrorenen Zwei-Fall-Audit
% (BCM_AUDIT_SYNTHESE_V1). T2-BIB integriert 2026-08-16 (36 verifizierte
% Einträge, 0 undefinierte Zitationen; prasad-Seiten am Primärdokument auf
% 91–114 korrigiert). Offen vor Upload: Retest-Widerleger, DE-Übersetzung,
% Dry-Run, Supplement der Zeilen-Ankertabellen (P1-5). Kein Kombi
% (Konvention 2026-08-15).
\documentclass[11pt,a4paper]{article}
\usepackage{cmap}
\usepackage[utf8]{inputenc}
\usepackage[T1]{fontenc}
\usepackage[ngerman]{babel}
\usepackage{lmodern}
\usepackage{microtype}
\usepackage{amsmath,amssymb,amsthm}
\usepackage{booktabs}
\usepackage{tabularx}
\usepackage{tikz}
\usetikzlibrary{arrows.meta,positioning,shapes.geometric,fit,calc}
% DE-Satzanpassung: Der deutsche Text ist laenger als der englische, wodurch das
% v1-F-Ledger von Abschnitt 2.3 ueber eine Seite hinauswaechst. ltablex macht
% tabularx seitenumbruchfaehig (Inhalt identisch, nur Umbruch erlaubt);
% emergencystretch faengt die langen deutschen Komposita in Ueberschriften ab.
\usepackage{ltablex}
\keepXColumns
\setlength{\emergencystretch}{2em}
\usepackage[colorlinks=true,linkcolor=blue,citecolor=blue,urlcolor=blue]{hyperref}
\usepackage{xurl}
\hypersetup{
  pdfencoding=unicode,
  psdextra,
  breaklinks=true,
  pdftitle={Das Bridge-Contract-Kriterium: Ein Form-Audit für Transportargumente, mit einer Instanziierung in der abc-Literatur},
  pdfauthor={Lukas Geiger},
  pdfsubject={Beweisauditierung, Bridge Contract, Transportargumente, abc-Vermutung},
  pdfkeywords={abc-Vermutung, Beweisauditierung, Transportargumente, Bridge Contract, Ledger-Audit, Mochizuki, Joshi}
}
\renewcommand*{\HyperDestNameFilter}[1]{de.#1}
\input{glyphtounicode}
\pdfgentounicode=1

\newtheorem{definition}{Definition}[section]
\newtheorem{lemma}{Lemma}[section]
\newtheorem{remark}{Bemerkung}[section]

\title{Das Bridge-Contract-Kriterium:\\Ein Form-Audit für Transportargumente,\\mit einer Instanziierung in der abc-Literatur}
\author{Lukas Geiger\\\small Unabhängiger Forscher, Bernau, Deutschland\\\small ORCID: 0009-0005-7296-1534}
\date{Oktober 2026 --- Version 0.3}

\begin{document}
\maketitle

\begin{abstract}
Manche Beweisbehauptungen widersetzen sich einer zeilenweisen Prüfung nicht
deshalb, weil einzelne Schritte undurchsichtig wären, sondern weil ihr
tragender Kern ein \emph{Transport} ist: Eine Größe wandert zwischen
theoretischen Rahmenwerken, und der numerische Gehalt der Behauptung ruht auf
der Zulässigkeit dieser Wanderung. Eine verwandte Klasse ruht auf
\emph{global-only}-Aussagen: Behauptungen, die erst nach Summation über alle
Stellen verfügbar sind, ohne stellenweisen Zugriff. Wir schlagen einen
siebenkomponentigen \emph{Bridge Contract} (Brückenvertrag) und einen
versionierten Ledger-Audit vor. Wir berichten eine retrospektive
Zwei-Fall-Machbarkeitsillustration des eingefrorenen
v1-T-Transport-Ledgers und eine substanzielle
\emph{selected-rows}-Instanziierung des Kriteriums in der abc-Literatur (die
dort bespielten v1-F-Zeilen werden benannt; die Vollmatrix-Kodierung aller
Zeilen ist zurückgestellt). Die Arbeit enthält kein RH-seitiges Verdikt;
RH-Uniformitäts- und global-only-Positivitätsbehauptungen bleiben als
ausdrücklich ungetesteter Ausblick stehen. Das Kriterium beurteilt
\emph{Vertragsform}, nicht mathematische Wahrheit: Ein bedientes Ledger ist
notwendig, nicht hinreichend, und der Auditor prüft das Vorhandensein und die
Adressierbarkeit einer Pflicht, niemals deren mathematische Gültigkeit. Sein
Wert liegt darin, dass sich Fehlschläge lokalisieren lassen -- der Audit benennt die
fehlende Klausel an einer angegebenen Zeile statt eines diffusen Zweifels.
Der aktuelle Forschungsstand gleicht die bedingte C7-Illustration an alle
fünf Hypothesen der unveränderten Vertragsdefinition an, nach einem
Gegenmodell zum früheren unvollständigen Diagramm. Spätere Entwicklungen
in Teilen 4--6 bleiben von den historischen Auditbindungen getrennt:
Gleichfeldrige Starrheit erlaubt globale Skalierung, während Koordinaten-,
Feld-, Perioden- und Volumenbrücken offen bleiben. Die Zwei-Fall-Rohmatrix
bleibt eine versionsgebundene Illustration, keine Reliabilitätsschätzung.
\end{abstract}

\clearpage
% sloppy nur fuer das Inhaltsverzeichnis: die langen deutschen Kapiteltitel
% ragen sonst in den rechten Rand.
\begingroup\sloppy
\tableofcontents
\endgroup
\clearpage

%=====================================================================
\section{Einleitung}\label{sec:intro}

Zwei Klassen von Beweisbehauptungen widersetzen sich der klassischen
zeilenweisen Prüfung auf charakteristische Weise. Die erste sind
\emph{Transportbehauptungen}: Eine Größe wird in einem Rahmenwerk konstruiert,
und von ihrer numerischen Ablesung wird behauptet, sie lande in einer
Zielregion, die in einem anderen Rahmenwerk festgelegt ist, wobei der Gehalt
der Behauptung auf der Zulässigkeit der Wanderung ruht und nicht auf
irgendeinem einzelnen Rechenschritt. Der dokumentierte Präzedenzfall ist der
umstrittene Transportschritt von Mochizukis Korollar 3.12, dessen
zehnjährige Kontroverse nach unserer Lesart nicht aus Uneinigkeit über
einzelne Zeilen besteht, sondern aus Uneinigkeit darüber, was der Transport
\emph{schuldet}~\cite{part1}. Die zweite sind
\emph{global-only-Behauptungen}: Aussagen, die erst nach Summation über alle
Stellen verfügbar sind, ohne stellenweisen Zugriff -- eine Gestalt, die
sowohl in der abc-Literatur als auch in Positivitätszugängen zur
Riemannschen Vermutung auftritt.

Diese Arbeit schlägt ein prüfbares Formkriterium für beide Klassen vor: den
\emph{Bridge Contract}, eine Liste von sieben Vertragskomponenten, die ein
Transport bedienen muss, bevor seine numerische Schlussfolgerung gezogen
werden darf, zusammen mit einem Ledger-artigen Auditprotokoll mit fester
Verdikt-Grammatik. Das Kriterium erwuchs aus den Audits dieser Serie -- der
Lückenlokalisierung von Teil 1~\cite{part1} und dem Zertifikatsinstrument
von Teil 3~\cite{part3} -- und ist seither, in den späteren Teilen, auf das
wichtigste publizierte Reparaturprogramm angewandt
worden~\cite{part4,part5,part6}. Hier stellen wir das Kriterium selbst als
übertragbare Methode vor: seine Definition
(Abschnitt~\ref{sec:criterion}), seine Stellung unter bestehenden
Auditierungsstrukturen (Abschnitt~\ref{sec:related}), eine
protokoll-eingefrorene retrospektive Zwei-Fall-Illustration des
eingefrorenen Transport-Ledgers (Abschnitt~\ref{sec:calibration}) und eine
\emph{selected-rows}-Instanziierung des Kriteriums am
abc-Transportschritt (Abschnitt~\ref{sec:abc}). Das begleitende
Landschaftspapier des Programms verzeichnet, wie sich das Kriterium unter
die von ihm auditierten Wege einordnet~\cite{landscapeA}.

Zwei Grenzen des Geltungsbereichs werden hier ausgesprochen statt der
Folgerung überlassen. Erstens ist das Instrument \emph{versioniert}: Das
volle Ledger von Abschnitt~\ref{sec:ledger} und das sechszeilige
Transport-Ledger der retrospektiven Illustration sind zwei verschiedene
Instrumente mit einem publizierten Crosswalk, und identische numerische
Zeilensuffixe bedeuten nicht identische Semantik. Zweitens berichtet diese
Arbeit genau eine Instanziierung, in der abc-Literatur, und diese
Instanziierung ist eine \emph{selected-rows}-Instanziierung: Die bespielten
Zeilen werden in Abschnitt~\ref{sec:abc} benannt, und eine Kodierung der
vollständigen v1-F-Matrix für diese Fälle ist auf eine spätere Version
zurückgestellt. Unter einer \emph{Instanziierung} verstehen wir durchgehend
die Anwendung des Kriteriums auf einen \emph{dispute complex}
(Streitkomplex) -- hier den abc-Transportstreit, der das ursprüngliche
Argument und sein Reparaturprogramm als einen einzigen Komplex mit zwei
auditierten Profilen umfasst. Uniformitätsbehauptungen in Richtung der
Riemannschen Vermutung und global-only-Positivitätsbehauptungen erscheinen
als ausdrücklich ungetesteter Ausblick (Abschnitte~\ref{sec:rh}
und~\ref{sec:global}); nirgends unten wird ein RH-seitiges Verdikt gezogen.

Damit die Zählung eindeutig bleibt, werden die in dieser Arbeit behandelten
Fälle durchgehend nach Art und nicht nach Ordnungszahl benannt: ein
\emph{retrospektiver Bekannte-Lücke-Fall} und eine \emph{akzeptierte
Kontrolle} (Abschnitt~\ref{sec:calibration}), die \emph{abc-Instanziierung}
mit einem \emph{vergleichenden Diskriminierungsfall} und einem
\emph{nachgelagerten Vererbungsbeispiel} (Abschnitt~\ref{sec:abc}), ein
\emph{externer Vergleich} (Abschnitt~\ref{sec:abc-external}) und der
\emph{RH-/global-only-Ausblick} (Abschnitte~\ref{sec:rh}
und~\ref{sec:global}).

Durchgehend wird eine Grenze ausnahmslos eingehalten: Das Kriterium
beurteilt \emph{Vertragsform}, nicht mathematische Wahrheit. Ein vollständig
bedientes Ledger ist notwendig für die Zulässigkeit einer Transportlesart,
niemals hinreichend für die zugrunde liegende Behauptung; eine unbediente
Zeile lokalisiert eine Pflicht, sie widerlegt kein Theorem. Der Auditor
prüft, ob eine Pflicht in einer benannten Version eines Textes vorhanden und
adressierbar ist, niemals, ob die Mathematik, die sie benennt, korrekt ist.
Der Wert, den wir für die Methode beanspruchen, ist genau diese
Lokalisierung: Audits, die in einer benannten fehlenden Klausel enden statt
in einem diffusen Zweifel.

%=====================================================================
\section{Das Kriterium}\label{sec:criterion}

\subsection{Der Bridge Contract}\label{sec:contract}

Die Situation, die wir axiomatisieren, ist die folgende. Ein Beweis behauptet
eine numerische Schlussfolgerung der Form \glqq die Ablesung eines im Rahmenwerk
$A$ konstruierten Objekts landet in einer im Rahmenwerk $B$ festgelegten
Zielregion\grqq{}, wobei die beiden Rahmenwerke konstruktionsbedingt keine
gemeinsame Auswertung teilen. Die Behauptung wird durch einen
\emph{Transport} vermittelt: Etiketten, Container oder Funktoren, die die
beiden Seiten in Beziehung setzen. Das Kriterium fragt, was der Transport
\emph{vertraglich} bereitstellen muss, bevor die numerische
Schlussfolgerung gezogen werden darf.

\begin{definition}[Bridge Contract]\label{def:bridge}
Eine Behauptung setze ein Quellobjekt zu einer numerischen Zielregion in
Beziehung. Ein \emph{Bridge Contract} $\mathrm{Br}$ für diese Behauptung
besteht aus den folgenden sieben Komponenten (im Folgenden
\emph{Komponente~1} bis \emph{Komponente~7}).
\begin{enumerate}
  \item \textbf{Quellobjekt mit Ablesung} ($X_q$, $\nu_q$, $i_q$): der
    quellseitige Bereich $X_q$, ein konkretes numerisches Ablesefunktional
    $\nu_q\colon X_q \to V$ in den geordneten Werteraum von Komponente~2 und
    die typisierte Abbildung $i_q\colon X_q \to C$ in den Container von
    Komponente~3 (im motivierenden Fall: eine log-volumetrische Ablesung des
    $q$-Piloten).
  \item \textbf{Zielobjekt mit Ziel-Ablesung und festgelegter Zielregion}
    ($X_\Theta$, $\nu_\Theta$, $(V,\preceq)$, $H_\Theta$): der zielseitige
    Bereich $X_\Theta$, eine Ziel-Ablesung $\nu_\Theta\colon X_\Theta \to V$,
    ein von beiden Ablesungen geteilter geordneter Werteraum $(V,\preceq)$
    und eine \emph{im Voraus} festgelegte Zielregion $H_\Theta \subseteq V$
    -- eine Abwärtsmenge $H_\Theta = \{v \in V : v \preceq \tau_\Theta\}$
    oder die entsprechende Aufwärtsmenge, mit der Schwelle
    $\tau_\Theta \in V$, die benannt wird, bevor der Quellwert inspiziert
    wird. Die typisierte Abbildung $i_\Theta\colon X_\Theta \to C$ gehört zu
    Komponente~3. Im motivierenden Fall ist $V = \mathbb{R}$ mit seiner
    üblichen Ordnung und $H_\Theta = \mathbb{R}_{\le -|\log\Theta|}$.
  \item \textbf{Typisierter gemeinsamer Container} ($C$, $\preceq_C$,
    $i_q\colon X_q \to C$, $i_\Theta\colon X_\Theta \to C$): ein gemeinsames
    Gefäß, ausgestattet mit der von der Behauptung verwendeten
    Ordnungsrelation $\preceq_C$ (eine Halbordnung oder die Präordnung, die
    die Quelle selbst angibt, mit ihrem Träger angegeben), zusammen mit
    \emph{explizit typisierten} Abbildungen von beiden Seiten, jede mit
    Definitions- und Zielbereich angegeben. Ein Container, in dem beide
    Objekte lediglich nebeneinander liegen, trägt für sich genommen keine
    Skala, und ein Container ohne angegebene Ordnung trägt keine
    Ordnungsbehauptung.
  \item \textbf{Zulässige Verknüpfung} ($L$): Etiketten- oder
    Verknüpfungsdaten, die eine \emph{erlaubte Identifikation} angeben, nicht
    bloße Homonymie. Einen Namen über Rahmenwerke hinweg zu teilen ist keine
    mathematische Relation.
  \item \textbf{Getrenntes Unbestimmtheitsbudget} ($I$): ein explizites
    Budget, das absorbierbare Renormierung, Verschmierung und
    \emph{nicht absorbierte Restlast} unterscheidet. Eine Unbestimmtheit, die
    aufgerufen wird, um eine Obstruktion zu absorbieren, muss gegen dieses
    Budget gebucht werden, nicht gegen den Container.
  \item \textbf{Kompatible gemeinsame Auswertung} ($\nu_C$). Eine typisierte
    Abbildung $\nu_C\colon C\to(V,\preceq)$ wird zusammen mit den von der
    Behauptung verlangten Kompatibilitätsidentitäten bereitgestellt,
    insbesondere $\nu_C\circ i_q=\nu_q$ auf dem angegebenen Quellbereich und,
    wo das Ziel durch eine Ziel-Ablesung $\nu_\Theta$ repräsentiert wird,
    $\nu_C\circ i_\Theta=\nu_\Theta$ auf dem angegebenen Zielbereich. Diese
    Kompatibilitätsklauseln sind verbindlich; die bloße Existenz einer
    numerischen Abbildung auf $C$ löst diese Komponente niemals ein.
  \item \textbf{Adressiertes Ordnungslemma.} Ein benanntes Lemma gibt die
    Hypothesen an, unter denen die zulässige Verknüpfung eine
    Ordnungsrelation in $C$ liefert, gibt die Monotonie- und
    Zielregionshypothesen an und leitet die verlangte Zielzugehörigkeit her.
    Für eine untere Zielmenge $H_\Theta=\{v\in V:v\preceq\tau_\Theta\}$
    lautet eine typische Form: Wenn (H1) $L(x,y)$ zulässig ist, (H2)
    $i_q(x)\preceq_C i_\Theta(y)$ gilt, (H3) $\nu_C$ auf solchen verknüpften
    Paaren monoton ist, (H4) die Kompatibilitätsklauseln von Komponente~6
    gelten und (H5) $\nu_C(i_\Theta(y))\preceq\tau_\Theta$, dann ist
    $\nu_q(x)\in H_\Theta$. Hypothese (H2) ist niemals eine freie Annahme:
    Der Vertrag verlangt vom Text, sie entweder in einem adressierten
    Teillemma aus den Verknüpfungsdaten von Komponente~4 \emph{herzuleiten}
    oder sie gesondert zu beweisen -- in beiden Fällen mit eigener
    Beweisadresse. Dieselbe Verankerungsforderung gilt für jede tragende
    Hypothese des Lemmas.
\end{enumerate}
\end{definition}

Ein Vertrags-Audit verzeichnet diese Komponente nur dann als textlich
bedient, wenn das Lemma, seine Hypothesen, seine behauptete Anwendbarkeit am
auditierten Schritt und eine Beweisadresse in der benannten Version vorhanden
und einsehbar sind. Der Auditor prüft Vorhandensein und Adressierbarkeit,
nicht mathematische Gültigkeit. Eine bediente Zeile zertifiziert daher nicht,
dass das Lemma oder sein Beweis korrekt ist.

Komponente~6 und Komponente~7 sind die tragenden Klauseln: Ohne sie liefert das
gemeinsame Nebeneinander in einem Container zusammen mit geteilten Etiketten
nur qualitative Nähe. In den unten auditierten Präzedenzfällen fallen dort
auch die schwersten Befunde an -- aber es sind nicht die einzigen Zeilen,
die Befunde tragen, und das Ledger berichtet das volle Profil statt einer
einzelnen übrig bleibenden Zeile.

Abbildung~\ref{fig:bridge_contract_architecture} illustriert die konzeptionelle Architektur eines Bridge Contracts und bildet das Zusammenspiel der sieben kanonischen Komponenten zwischen den disjunkten Quell- und Ziel-Frameworks ab. Tabelle~\ref{tab:vertrags_komponenten} fasst die sieben kanonischen Komponenten, ihre mathematischen Signaturen, die geforderten Verifikationspflichten und die charakteristischen Fehlermodi synoptisch zusammen.

\begin{figure}[p]
\centering
\resizebox{0.88\textwidth}{!}{%
\begin{tikzpicture}[
  >=Stealth,
  box/.style={draw=black!70, rounded corners=3pt, line width=0.8pt, align=center, inner sep=4pt},
  srcbox/.style={box, fill=blue!6, draw=blue!70!black},
  tgtbox/.style={box, fill=orange!8, draw=orange!80!black},
  ctrbox/.style={box, fill=teal!8, draw=teal!80!black},
  valbox/.style={box, fill=purple!6, draw=purple!70!black},
  gatebox/.style={box, fill=gray!8, draw=black!60, dashed},
  arrow/.style={->, line width=0.9pt, draw=black!80},
  annot/.style={font=\scriptsize, text=black!70, align=left}
]

  % Top layer: Source, Container, Target
  \node[srcbox, text width=3.2cm] (src) at (0, 7.0) {
    \textbf{Framework $\mathcal{F}_A$ (Quellseite)}\\[2pt]
    \footnotesize Quellbereich $X_q$\\[1pt]
    \scriptsize Ablesung $\nu_q\colon X_q \to V$\\
    \textit{(Komponente~1)}
  };

  \node[ctrbox, text width=4.4cm] (ctr) at (4.7, 7.0) {
    \textbf{Gemeinsamer Behälter $(C, \preceq_C)$}\\[2pt]
    \footnotesize Typisiertes Gefäß mit Ordnung $\preceq_C$\\[1pt]
    \scriptsize Abbildungen $i_q\colon X_q \to C$, $i_\Theta\colon X_\Theta \to C$\\
    \textit{(Komponente~3)}
  };

  \node[tgtbox, text width=3.2cm] (tgt) at (9.4, 7.0) {
    \textbf{Framework $\mathcal{F}_B$ (Zielseite)}\\[2pt]
    \footnotesize Zielbereich $X_\Theta$\\[1pt]
    \scriptsize Ziel-Ablesung $\nu_\Theta\colon X_\Theta \to V$\\
    \textit{(Komponente~2)}
  };

  % Intermediate layer: Linkage and Budget
  \node[gatebox, text width=3.3cm] (link) at (2.2, 3.8) {
    \textbf{Zulässige Verknüpfung $L$}\\[1pt]
    \scriptsize Erlaubte Identifikation (keine Homonymie)\\\textit{(Komponente~4)}
  };

  \node[gatebox, text width=3.4cm] (bdg) at (7.2, 3.8) {
    \textbf{Unbestimmtheits-}\\\textbf{budget $I$}\\[1pt]
    \scriptsize Absorbierte vs.\ verbleibende Last\\\textit{(Komponente~5)}
  };

  % Value Space
  \node[valbox, text width=7.2cm] (val) at (4.7, -1.0) {
    \textbf{Geordneter Werteraum $(V,\preceq)$}\\[2pt]
    \footnotesize Vorab fixierter Zielbereich $H_\Theta = \{v \in V : v \preceq \tau_\Theta\}$\\[1pt]
    \scriptsize Feste Schranke $\tau_\Theta \in V$ vorab benannt
  };

  % Embeddings
  \draw[arrow] (src.east) -- (ctr.west) node[midway, above, font=\scriptsize] {$i_q$};
  \draw[arrow] (tgt.west) -- (ctr.east) node[midway, above, font=\scriptsize] {$i_\Theta$};

  % Downward evaluation
  \draw[arrow, line width=1.1pt, draw=teal!80!black] (ctr.south) -- (val.north)
    node[pos=0.85, left, font=\scriptsize] {$\nu_C\colon C \to (V,\preceq)$};

  \node[annot, text width=3.8cm, align=center] at (7.2, 1.3) {
    \textbf{Komponente~6}:\\
    $\nu_C \circ i_q = \nu_q$\\
    $\nu_C \circ i_\Theta = \nu_\Theta$
  };

  % Direct readouts to value space
  \draw[arrow, dashed, draw=blue!70!black] (src.south) to[out=-90,in=165]
    node[left, font=\scriptsize] {$\nu_q$} (val.north west);
  \draw[arrow, dashed, draw=orange!80!black] (tgt.south) to[out=-90,in=15]
    node[right, font=\scriptsize] {$\nu_\Theta$} (val.north east);

  % C7: north anchor below the value node; extra hypotheses grow down.
  \node[draw=black!70, rounded corners=4pt, line width=0.8pt, fill=yellow!8,
    text width=11.2cm, align=center, inner sep=5pt, anchor=north]
    (lem) at ($(val.south)+(0,-0.55)$) {
    \textbf{Komponente~7: Bedingtes adressiertes Ordnungslemma}\\[2pt]
    \scriptsize (H1) $L(x,y)$ zulässig\\
    \scriptsize (H2) $i_q(x)\preceq_C i_\Theta(y)$\\
    \scriptsize (H3) $\nu_C$ auf dem verknüpften Paar monoton\\
    \scriptsize (H4) $\nu_C\circ i_q=\nu_q$, $\nu_C\circ i_\Theta=\nu_\Theta$ auf den angegebenen Bereichen\\
    \scriptsize (H5) $\nu_C(i_\Theta(y))\preceq\tau_\Theta$\\[2pt]
    \scriptsize $(\mathrm{H1})\land\cdots\land(\mathrm{H5})\ \Longrightarrow\ \nu_q(x)\in H_\Theta$
  };

  \draw[arrow, draw=black!60, line width=0.7pt] (val.south) -- (lem.north);

\end{tikzpicture}%
}
\caption{Bedingte Bridge-Contract-Illustration für die untere Zielmenge.
Alle fünf Hypothesen werden benötigt; (H2) und jede tragende Pflicht
verlangen eigene Beweisadressen. Das Diagramm illustriert den Vertrag
und zertifiziert nicht die Korrektheit eines Quellenbeweises.}
\label{fig:bridge_contract_architecture}
\end{figure}

\paragraph{Entwicklung der C7-Illustration.}
Das frühere Diagramm ließ die Zielschranke (H5) weg und ersetzte
Kompatibilität (H4) durch eine beschränkte Restlast. Seine Implikation
scheitert im Identitätsmodell: Alle Bereiche und Träger seien $\mathbb R$,
alle dargestellten Abbildungen Identitäten, mit üblicher Ordnung,
$x=y=1$, zulässigem $L(x,y)$, Restlast $0$, $\tau_\Theta=0$ und
$H_\Theta=(-\infty,0]$. Die früher dargestellten Prämissen gelten,
einschließlich Kompatibilität, aber $\nu_q(x)=1\notin H_\Theta$.
Dies widerlegt die unvollständige Bildimplikation, nicht
Definition~\ref{def:bridge} mit (H1)--(H5).
Abbildung~\ref{fig:bridge_contract_architecture} stellt diese fünf
Bedingungen nun getrennt dar; die Illustration zertifiziert keinen
Quellenbeweis und begründet keine experimentelle Validierung der Methode.



\begin{table}[htbp]
\centering
\footnotesize
\renewcommand{\arraystretch}{0.95}
\caption{Die sieben kanonischen Komponenten eines Bridge Contracts $\mathrm{Br}$ und ihre Verifikationspflichten.}
\label{tab:vertrags_komponenten}\label{tab:contract_components}
\begin{tabularx}{\textwidth}{@{} >{\raggedright\arraybackslash}p{0.18\textwidth} >{\raggedright\arraybackslash}p{0.22\textwidth} >{\raggedright\arraybackslash}X >{\raggedright\arraybackslash}X@{}}
\toprule
Komponente & Mathematische Signatur & Verifikationspflicht & Fehlermodus / Lückentypologie \\
\midrule
C1: Quell-Ablesung & $(X_q, \nu_q, i_q)$ & Expliziter Definitions-, Zielbereich und Auswertungsfunktional $\nu_q\colon X_q \to V$ & Untypisierte Quelle, implizite Ablesung \\
\addlinespace[3pt]
C2: Zielregion & $(X_\Theta, \nu_\Theta, (V,\preceq), H_\Theta)$ & Werteraum $(V,\preceq)$ und Zielregion $H_\Theta$ \emph{im Voraus} fixiert & Verschobene Schwelle, nachträgliche Zielmenge \\
\addlinespace[3pt]
C3: Typisierter Container & $(C, \preceq_C, i_q, i_\Theta)$ & Träger $C$ mit Halbordnung/Präordnung $\preceq_C$ und typisierten Abbildungen & Skalenloses Nebeneinander, ungeordneter Träger \\
\addlinespace[3pt]
C4: Zulässige Verknüpfung & $L(x,y)$ & Erlaubte Identifikationsregel durch Satz/Hypothese spezifiziert & Nominale Identität, bloße Homonymie \\
\addlinespace[3pt]
C5: Unbestimmtheitsbudget & $(I, \mathrm{Renorm},\allowbreak \mathrm{Blur},\allowbreak \mathrm{Res})$ & Getrennte Verbuchung herausgeteilter Gruppenoperationen vs.\ Restlast & Rest dem Container angelastet, ungebuchter Drift \\
\addlinespace[3pt]
C6: Kompatible Auswertung & $\nu_C\colon C \to V$ & Kom\-mu\-ta\-ti\-vi\-täts\-iden\-ti\-tä\-ten $\nu_C\circ i_q=\nu_q$ und $\nu_C\circ i_\Theta=\nu_\Theta$ bewiesen & Koinzidenz unterstellt, fehlende Kommutation \\
\addlinespace[3pt]
C7: Adressiertes Ordnungslemma & Lemma $\mathrm{Ord}(\dots)$, (H1)--(H5) & Monotonie von $\nu_C$, Herleitung von (H2) $i_q(x)\preceq_C i_\Theta(y)$ & \glqq Follows formally\grqq{}, unadressierter Ordnungssprung \\
\bottomrule
\end{tabularx}
\end{table}

\subsection{Zur Lesart der Bereichsausdrücke}\label{sec:glossary}

Vier Ausdrücke in Definition~\ref{def:bridge} sind generisch gemeint und
werden hier fixiert, damit der Vertrag auch außerhalb des motivierenden Falls
instanziiert werden kann. Eine Identifikation ist \emph{erlaubt}, wenn sie
durch eine benannte Hypothese oder eine benannte Klasse von Theoremen
zugelassen wird -- geteilte Notation ist keine Lizenz. Eine Unbestimmtheit
ist \emph{absorbierbar}, wenn sie durch eine explizit angegebene Äquivalenz-
oder Invarianzrelation herausgeteilt werden kann. Die \emph{Restlast} ist der
Term, der übrig bleibt, nachdem jeder erlaubte Quotient gebildet wurde; sie
wird gegen das Budget von Komponente~5 gebucht und niemals gegen den
Container. Ein \emph{zweiter Modus} ist ein alternativer Zweig, eine
alternative Normierung oder Vorzeichenkonvention, die den Zielwert oder die
Ordnungsrelation ändern kann und deren Vorhandensein die entsprechende
Ledger-Zeile auslöst. Das inter-universelle Teichmüller-Vokabular des
motivierenden Falls liefert Beispiele für diese Ausdrücke, nicht ihre
Definitionen.

\subsection{Ledger-Form und Verdikt-Grammatik}\label{sec:ledger}

\noindent\textbf{Instrumentversionen.} Wir unterscheiden zwei
versionsskopierte Instrumente. \textbf{BCM-Ledger v1-F} ist das volle Ledger,
mit dem das allgemeine siebenkomponentige Kriterium in dieser Arbeit
operationalisiert wird. \textbf{BCM-Ledger v1-T} ist das sechszeilige
Transportinstrument, das am 15.~August 2026 vor den beiden historischen
Quellenkodierungsdurchgängen eingefroren wurde. Die retrospektive
Zwei-Fall-Illustration in Abschnitt~\ref{sec:calibration} betrifft
ausschließlich v1-T. Zeilenbezeichner sind durch das Versionspräfix
skopiert; identische numerische Suffixe bedeuten nicht identische Semantik.
Die Zuordnungstabelle unten ist ein Crosswalk, keine Behauptung, dass die
beiden Instrumente austauschbar wären.

\medskip

Der Vertrag wird durch ein festes Ledger von Testzeilen auditiert. Jede Zeile
erhält eines von drei Verdikten -- \emph{liefert} / \emph{liefert nicht} /
\emph{liefert teilweise} -- zusammen mit einem Konfidenzqualifikator und
einem Quellenanker (Seiten- bzw.\ Abschnittsangabe des auditierten Textes).
Mehrere Zeilen sind aufgespalten, eine aus den Audits selbst gelernte
Lektion: Ohne die Aufspaltung fallen zwei verschiedene Befunde in ein Wort
zusammen. Tabelle~\ref{tab:bcm_ledger_v1f} ist normativ für v1-F. \glqq Primärkomponente\grqq{}
benennt die lokalisierte Vertragspflicht; Abhängigkeiten werden transparent
gelistet.

\begin{table}[htbp]
\centering
\footnotesize
\renewcommand{\arraystretch}{0.70}
\setlength{\tabcolsep}{3.5pt}
\caption{Das vollständige BCM-Ledger v1-F: Kern-Prüfzeilen, primäre Vertragskomponenten und beobachtbare Prüfkriterien.}
\label{tab:bcm_ledger_v1f}
\begin{tabularx}{\textwidth}{@{}l l >{\raggedright\arraybackslash}X@{}}
\toprule
v1-F-Zeile & Primär & Beobachtbarer Test der Bedienung \\
\midrule
F-R0 & C1 (abh.\ C2)
  & Quelle, Ziel, Ablesungen, Werteraum und Zielregion sind getrennt mit
    Definitions- und Zielbereich typisiert. \\
F-R1 & C3
  & Ein gemeinsamer Container und beide typisierten Einbettungen oder
    Abbildungen sind benannt. \\
F-R2a & C4 (abh.\ C6)
  & Für jeden Schritt innerhalb des Links ist das Transportgesetz mit einer
    Adresse und einer sichtbaren Wirkung auf die Ablesung angegeben;
    \glqq keine Operation\grqq{} ergibt \texttt{X}, nicht automatisch \texttt{E+}. \\
F-R2b & C6 (abh.\ C3)
  & Der containerübergreifende Schritt trägt einen sichtbaren
    Umrechnungsfaktor oder eine Kompatibilitätsgleichung; bloße Konvention
    oder Koinzidenz ist höchstens \texttt{I}. \\
F-R3 & C2
  & $H_\Theta$ bzw.\ $\tau_\Theta$ wird festgelegt, bevor der Quellwert
    inspiziert wird; kein nachträgliches Verschieben. \\
F-R4a & C5
  & Jede nicht neutrale Last oder Kompensation ist quantifiziert und einer
    Budgetkategorie zugeordnet. \\
F-R4b & C5
  & Die Last wird am Brückenschritt selbst gebucht, mit Träger, Vorzeichen
    und Restterm; nachgelagert zu einer vorausgesetzten Schlussfolgerung
    genügt nicht. \\
F-R5 & C5
  & Absorbierbare Renormierung, Verschmierung und nicht absorbierter Rest
    werden auseinandergehalten; der Rest wird nicht dem Container
    angelastet. \\
F-R6a & C6
  & Die gemeinsame Auswertung und alle von der Behauptung benötigten
    Kompatibilitätsidentitäten sind im Text vorhanden und adressiert. \\
F-R6b & C7
  & Das benannte Ordnungslemma listet seine Hypothesen auf, behauptet ihre
    Anwendbarkeit am Schritt und trägt eine konkrete Beweisadresse;
    zusätzlich trägt jede tragende Hypothese -- insbesondere (H2) --
    ihren eigenen Anker (Herleitung aus den Verknüpfungsdaten oder
    gesonderter Beweis). Der Auditor prüft Vorhandensein und
    Adressierbarkeit dieser Anker, niemals ihre Gültigkeit. \\
F-R7 & C7 (abh.\ C2)
  & Die tatsächlich hergeleitete Aussage ist die Zugehörigkeit zu
    $H_\Theta$; Approximation, Nähe oder asymptotische Annäherung wird nicht
    zu Zugehörigkeit umetikettiert. \\
\midrule
F-R8 & \emph{Erweiterungszeile} E8
  & \textbf{E8 (Definition).} E8 ist eine Bereichserweiterung außerhalb des
    siebenkomponentigen Crosswalks: Sie bildet auf keine Vertragskomponente
    ab und ist nur als ausgelöste Diagnostik Teil von v1-F.
    \emph{Auslöser:} Im Bereichsmodell der Behauptung sind mehrere Modi,
    Zweige, Normierungen oder ein vorzeichen- bzw.\ ordnungsrelevanter
    Alternativmodus vorhanden. Dann müssen alle Modi sichtbar und entweder
    vom Ordnungslemma abgedeckt oder mit Gründen ausgeschlossen sein (hier
    berührt E8 C5 und C7, ohne zu ihnen zu gehören). Ohne den Auslöser ist
    die Zeile notwendig \texttt{X}. \\
\bottomrule
\end{tabularx}
\end{table}

Jede aktive \emph{Kernzeile} (F-R0 bis F-R7) hat genau eine primäre
Vertragskomponente. Unterstützende Abhängigkeiten dürfen gelistet werden,
aber ein negatives Verdikt wird an der im Crosswalk benannten
Primärkomponente berichtet. Die Erweiterungszeile F-R8 ist die einzige
Ausnahme: Ihrer E8-Definition nach bildet sie auf keine Komponente ab, und
ein ausgelöster F-R8-Befund wird am Erweiterungsetikett E8 selbst berichtet,
niemals an einer Vertragskomponente. F-R8 ist keine universelle Kernzeile;
sie wird nur durch den angegebenen Mehrmoden-Auslöser aktiviert und ist
andernfalls als X kodiert.

Die Kodierung verwendet fünf Rohcodes: explizit eingelöst (\texttt{E+}),
explizit nicht einlösbar (\texttt{E-}), nur über eine dokumentierte
Rekonstruktion einlösbar (\texttt{I}), nicht adressiert (\texttt{N}) und
außerhalb des Geltungsbereichs (\texttt{X}), letzterer nur mit expliziter
Geltungsbereichs- oder Auslöserbegründung zulässig. \texttt{I} und \texttt{N}
sind Formbefunde, keine Fehlerbefunde. Rohcodes werden von der öffentlichen
Verdikt-Grammatik niemals gelöscht: E+ bildet auf \emph{liefert} ab; I bildet
auf \emph{liefert teilweise durch dokumentierte Rekonstruktion} ab; E$-$ und
N bilden auf \emph{liefert nicht} ab, wobei die Unterscheidung zwischen
expliziter Offenheit und Nichtadressierung erhalten bleibt; X ist außerhalb
des Geltungsbereichs. Kodiererpaare werden weiterhin als Paare berichtet.
Eine mit \texttt{X} kodierte Zeile trägt kein Verdikt und wird nicht im
Nenner gezählt. Wo eine Zeile zweifach kodiert wird, ist die in der
retrospektiven Illustration verwendete Doppelkodierungsregel in
Abschnitt~\ref{sec:calibration} angegeben; sie ist eine Regel jener
Instrumentversion und nicht Teil von v1-F.

\subsection{Crosswalk zwischen den beiden Instrumentversionen}\label{sec:crosswalk}

Das eingefrorene Transport-Ledger v1-T steht Teilen von v1-F funktional nahe,
ist mit ihnen aber nicht identisch. Die Zuordnung wird publiziert, damit die
retrospektive Illustration von Abschnitt~\ref{sec:calibration} gegen das
volle Instrument gelesen werden kann, ohne dass eines durch das andere
ersetzt wird.

\begin{table}[htbp]
\centering
\caption{Crosswalk zwischen dem eingefrorenen Transportinstrument (BCM-Ledger v1-T) und dem Vollinstrument (BCM-Ledger v1-F).}
\label{tab:crosswalk_v1t_v1f}
\begin{tabularx}{\textwidth}{@{}l >{\raggedright\arraybackslash}X l@{}}
\toprule
v1-T-Zeile & Eingefrorene Semantik & Nächstes in v1-F \\
\midrule
T-R2a & Transportgesetz explizit & F-R2a / C4 \\
T-R2b & sichtbarer Umrechnungsfaktor & F-R2b und F-R6a / C6 \\
T-R4a & Kompensationsmechanismus benannt und lokalisiert & F-R4a / C5 \\
T-R4b & Existenz und Wohldefiniertheit bewiesen & \emph{kein direktes
  Äquivalent}$^{\dagger}$ \\
T-R6 & zweite Modi und Fälle sichtbar und symmetrisch & F-R8 / ausgelöstes
  E8 \\
T-R8 & kettenschließende Kompatibilität hergeleitet, lokal oder global
  & F-R6a, F-R6b / C6--C7 \\
\bottomrule
\end{tabularx}
\end{table}

\noindent {\small $^{\dagger}$\,T-R4b hat kein direktes v1-F-Äquivalent: Der
nächste Nachbar ist F-R4b (Buchung der Last am Brückenschritt), aber die
eingefrorene Existenz- und Wohldefiniertheitssemantik von T-R4b wird selbst
von keiner v1-F-Kernzeile getragen. Der Crosswalk verzeichnet Nähe, nicht
Identität.}

\noindent
Tabelle~\ref{tab:crosswalk_v1t_v1f} zeigt funktionale Nähe, nicht Identität. Zu beachten ist
insbesondere, dass sich die numerischen Suffixe kreuzen: Die
Zweiter-Modus-Zeile von v1-T ist T-R6, während die Zweiter-Modus-Zeile von
v1-F F-R8 ist. Das ist der Grund, weshalb in dieser Arbeit an jedem
Zeilenbezeichner das Versionspräfix mitgeführt wird.

\subsection{Geltungsbereich des Kriteriums}\label{sec:scope}

Das Kriterium beurteilt \emph{Vertragsform}, nicht Wahrheit. Was ein
vollständig bedientes Ledger feststellt, ist textliche Vertragsvollständigkeit
für die Transportlesart eines ausgewählten Schrittes; es ist kein Beweis der
zugrunde liegenden Behauptung, und es zertifiziert nicht die Mathematik, die
die bedienten Zeilen benennen. Umgekehrt widerlegt eine unbediente Zeile die
Behauptung nicht -- sie lokalisiert die Last: Die Zeile benennt die Klausel,
die ein vollständiges Argument bedienen muss. Diese Asymmetrie ist der Zweck
des Kriteriums. Die Auditeinheit ist entsprechend eng: Das Ledger zerlegt
einen ausgewählten Transportstreit in benannte Pflichten; es reduziert nicht
die historische Kontroverse als Ganzes. Was ein Audit liefert, ist ein Profil
benannter Pflichten an einem ausgewählten umstrittenen Transportschritt, mit
Gewichten -- und dort kann dann mathematische Arbeit statt Exegese
investiert werden.

%=====================================================================
\section{Verwandte Auditierungsstrukturen}\label{sec:related}

Fünf Arbeitsfelder berühren das Kriterium von Abschnitt~\ref{sec:criterion}:
normative Darstellungen dazu, wann eine Auslassung in einem mathematischen
Beweis akzeptabel ist; maschinelle Formalisierung tiefer Mathematik;
Zertifikatsschnittstellen in der Informatik; strukturierte
Assurance-Argumentation in der Sicherheitstechnik; und die Methodologie
dokumentierter mathematischer Uneinigkeit. Zwei Beobachtungen ordnen das
Folgende. Erstens ist das nächste heuristische Analogon unter den hier
verglichenen Strukturen die Tradition der zertifizierenden
Algorithmen~\cite{McConnell2011}, von der sich das Kriterium in der Richtung
des Zertifikats und darin unterscheidet, was die Prüfung garantiert
(Abschnitt~\ref{sec:related-certificates}). Zweitens ist uns innerhalb des
hier gesichteten Korpus keine unabhängige, begutachtete Fallanalyse der
abc/IUT-Kontroverse bekannt -- eine Lücke, die fortbesteht und die die
vorliegende Arbeit nicht schließt, da ihr Beitrag auf der Ebene der
Vertragsform liegt (Abschnitte~\ref{sec:related-controversy}
und~\ref{sec:scope}).

\subsection{Beweisauditierung und Peer-Review-Methodologie}\label{sec:related-audit}

Andersen~\cite{Andersen2020} fragt, wann eine Lücke in einem publizierten
Beweis als akzeptabel behandelt werden darf, und gibt normative Kriterien für
die Substanz des Ausgelassenen. Der Bridge Contract beantwortet eine andere
Frage: ob die Vertragskomponenten von Abschnitt~\ref{sec:contract} in der
Darstellung überhaupt bedient werden. Die beiden Beurteilungen sind
unabhängig -- ein vollständig bedientes Ledger ist eine Aussage über Form,
nicht über die Wahrheit der transportierten Behauptung.
De~Millo--Lipton--Perlis~\cite{DeMillo1979} beschreiben Beweisakzeptanz als
sozialen Prozess, getragen von gemeinschaftlicher Lektüre. Das Ledger ist
unabhängig von der Aufnahme durch eine Gemeinschaft und von einem bestimmten
institutionellen Begutachtungsverfahren; auditorunabhängig ist es nicht, wie
der berichtete Kodierungsdissens zeigt
(Abschnitt~\ref{sec:calibration}). Andersen~\cite{Andersen2017} dokumentiert,
was Gutachter in der Mathematik tatsächlich prüfen -- eine deskriptive
Darstellung eines bestehenden Prozesses; das Kriterium ist ein Instrument,
das außerhalb jedes solchen Prozesses angewandt werden kann.
Avigad~\cite{Avigad2021} behandelt die Verlässlichkeit informeller
mathematischer Schlussfolgerungen im Allgemeinen; das Kriterium adressiert
einen Spezialfall innerhalb dieser Spannweite: einen Transportschritt, dessen
Schlussfolgerung eine numerische Zugehörigkeitsbehauptung ist.

\subsection{Formale Verifikation und Lean}\label{sec:related-formal}

Scholzes Liquid Tensor Experiment~\cite{Scholze2022LTE} unterwarf bewusst ein
umstrittenes eigenes Resultat einem maschinenprüfbaren Vertrag. Es beantwortet
die Wahrheitsfrage durch vollständige Formalisierung, zu hohen Kosten und
unter Beteiligung des Autors; das Ledger beantwortet die Formfrage am
publizierten Text und ohne Mitwirkung des Autors. Die mathematische Bibliothek
\emph{mathlib}~\cite{Mathlib2020} ist die Infrastruktur, in der solche
formalen Verträge formuliert und wiederverwendet werden; das Kriterium
arbeitet eine Ebene früher, an natürlichsprachlichen Argumenten, die nicht
formalisiert worden sind.
Buzzard--Commelin--Massot~\cite{Buzzard2020} dokumentieren, was die
Formalisierung moderner arithmetischer Geometrie in der Praxis kostet; dieser
Aufwand ist der Grund, weshalb ein leichtgewichtiges Formkriterium neben der
vollständigen Formalisierung einen Platz hat statt in Konkurrenz zu ihr.
Hales et al.~\cite{Hales2017} verifizierten die Keplersche Vermutung
maschinell, nachdem die konventionelle Begutachtung an ihre Grenze gestoßen
war; das Kriterium ist darauf angelegt, \emph{vor} dem Abschluss der
Begutachtung anwendbar zu sein, sodass die umstrittene Stelle zu dem
Zeitpunkt benennbar ist, zu dem der Streit entsteht.

\subsection{Zertifikatsschnittstellen in der Informatik}\label{sec:related-certificates}

Dieses Feld enthält das nächste heuristische Analogon unter den hier
verglichenen Strukturen.
Ein zertifizierender Algorithmus im Sinne von McConnell et
al.~\cite{McConnell2011} liefert zusammen mit seiner Ausgabe einen Zeugen,
den ein einfacher und unabhängig prüfbarer Verifizierer akzeptiert; die
Korrektheit des einzelnen Laufs wird dadurch festgestellt, ohne der
Implementierung zu vertrauen. Der Bridge Contract hat eine heuristische
Schnittstellenanalogie zu zertifizierenden Algorithmen: Beide stellen ein
einsehbares Artefakt neben eine Behauptung. Die Analogie endet vor der
Korrektheitszertifizierung. Die Unterschiede sind systematisch und nicht
zufällig; sie sind in Tabelle~\ref{tab:certifying_vs_bcm} zusammengefasst:

\begin{table}[htbp]
\centering
\small
\caption{Struktureller Vergleich zwischen zertifizierenden Algorithmen und Bridge-Contract-Audits.}
\label{tab:certifying_vs_bcm}
\begin{tabularx}{\textwidth}{@{}l >{\raggedright\arraybackslash}X >{\raggedright\arraybackslash}X@{}}
\toprule
Merkmal & Zertifizierender Algorithmus & Bridge-Contract-Audit \\
\midrule
Erzeuger des Artefakts & der Algorithmus oder Erzeuger & eine dritte Partei,
  im Nachhinein \\
Zeitpunkt der Ausstellung & mit der Ausgabe & nach Publikation des Textes \\
Prüfer & formal und einfach & fachlich und hermeneutisch \\
Entscheidbarkeit & typischerweise algorithmisch & nicht garantiert \\
Garantie & Korrektheit des Laufs oder der Ausgabe relativ zum Prüfer
  & nur textliche Vertragsvollständigkeit \\
Dissens & kein notwendiger Bestandteil & ausdrücklich möglich und zu
  berichten \\
\bottomrule
\end{tabularx}
\end{table}

\noindent
Neculas Proof-Carrying Code~\cite{Necula1997} setzt eine im Voraus
vereinbarte Policy voraus, gegen die der Erzeuger einen maschinenprüfbaren
Nachweis der Einhaltung liefert. Der in Abschnitt~\ref{sec:abc} behandelte
Präzedenzfall ist genau einer, in dem eine solche Policy nie vereinbart
wurde: Das Ledger rekonstruiert die fehlende Vereinbarung im Nachhinein und
benennt die Zeile, an der sie nie geschlossen wurde.
Blum--Kannan~\cite{Blum1995} prüfen Resultate zur Laufzeit, wiederholt und
billig, über einen Strom von Instanzen; das Kriterium wird einmal angewandt,
auf eine einzelne argumentative Stelle in einem Text.

\subsection{Assurance Cases und Goal Structuring Notation}\label{sec:related-assurance}

Rushby~\cite{Rushby2015} fragt, wann ein strukturiertes Argument trägt, und
liefert das Vokabular für die hier maßgeblichen Fehlermodi -- Defeater (Entkräftungsfaktoren) und
Bestätigungsfehler bei der Partei, die das Argument konstruiert --, wobei er
Überzeugungskraft in abgestuften, holistischen Begriffen beurteilt. Die
Verdikt-Grammatik von Abschnitt~\ref{sec:ledger} ist bewusst gröber und
lokaler: Jede Zeile erhält eines von drei Verdikten über eine
Vertragsklausel, wobei der Konfidenzqualifikator getrennt vom Verdikt selbst
festgehalten wird. Die Goal Structuring Notation (GSN)~\cite{GSN2021}
standardisiert, wie Behauptung, Argument und Evidenz auseinandergehalten
werden, und ist von der Anlage her domänenneutral; die sieben
Vertragskomponenten von Abschnitt~\ref{sec:contract} sind die komplementäre
Wahl, inhaltlich für eine wiederkehrende Situation bestimmt und daher nicht
ohne Anpassung übertragbar. Bloomfield--Rushby~\cite{Bloomfield2021} machen
die explizite Behandlung von Defeatern verbindlich, und zwar für einen
Zertifizierungskontext, in dem ein Argument den Einsatz autorisiert. Das
Kriterium trägt keine Zulassungsautorität dieser Art: Es ist ein
Leseinstrument, und seine Verdikte richten sich an Leser eines Beweises und
nicht an eine Aufsichtsbehörde.

\subsection{Fehleranalysen und Kontroversen-Methodologie}\label{sec:related-controversy}

Aberdein~\cite{Aberdein2023} erklärt, weshalb tiefe Uneinigkeiten in der
Mathematik selbst unter kompetenten Parteien fortbestehen, und arbeitet auf
der Ebene der Uneinigkeit selbst. Das Kriterium arbeitet eine Ebene darunter:
Es lokalisiert die einzelne strittige transportierte Zeile und bezieht zu ihr
keine Position. Das Paar Scholze--Stix~\cite{ScholzeStix2018} und
Mochizuki~\cite{Mochizuki2019} ist der Referenzfall des Feldes und ist nur
als Paar aufschlussreich -- ein stellengenauer Fehlerbericht und seine
Erwiderung, in denen die beiden Seiten mangels eines geteilten Formrasters
aneinander vorbeireden. Das Ledger von Abschnitt~\ref{sec:ledger} ist die
verallgemeinerte und übertragbare Fassung eines solchen Rasters.
Castelvecchi~\cite{Castelvecchi2020} dokumentiert die institutionellen Folgen
dieser Pattsituation, wie sie außerhalb des Feldes sichtbar wurden.
Joshi~\cite{Joshi2025Report} berichtet über die Kontroverse aus der Position
einer an ihr beteiligten Partei; die Arbeit wird hier als Parteiäußerung
zitiert und nicht als unabhängige Bestätigung irgendeiner Lesart,
einschließlich unserer eigenen.

Darüber hinaus haben wir in dem für diese Arbeit verwendeten begrenzten,
nicht systematischen Korpus keine unabhängige begutachtete Detailanalyse der
abc/IUT-Kontroverse identifiziert; dies ist kein Befund einer systematischen
Übersichtsarbeit. Was dieser Korpus enthält, ist Primärliteratur der
Parteien, journalistische Berichterstattung~\cite{Castelvecchi2020} und
allgemeine Theorie tiefer Uneinigkeit~\cite{Aberdein2023}. Wir verzeichnen
die Abwesenheit als fortbestehende Lücke und nicht als eine, die die
vorliegende Arbeit schließt, und wir beanspruchen nicht, erschöpfend gesucht
zu haben. Was das Kriterium beiträgt, steht dazu orthogonal: ein Formverdikt
über einen einzelnen Transportschritt, keine historische Darstellung der
Debatte (Abschnitt~\ref{sec:limits}).

\subsection{Stellung des Kriteriums}\label{sec:related-position}

Das Kriterium sitzt im Schnittpunkt dieser fünf Felder und nicht in der
Nachfolge eines einzelnen von ihnen. Es entlehnt die Schnittstelle des
einsehbaren Artefakts der Tradition zertifizierender Algorithmen ohne deren
Korrektheitsgarantie, die explizite Trennung von Behauptung, Argument und
Evidenz der Assurance-Case-Praxis, die Frage nach der zulässigen Auslassung
der Beweisauditierungsliteratur und die Lektüre eines umstrittenen
Transportschrittes der dokumentierten Kontroverse -- und wendet sie auf eine
Klasse von Argumenten an, für die vollständige Formalisierung im Prinzip
verfügbar ist, bei gegenwärtigen Kosten aber selten in der Praxis. Seine
Verdikte bleiben durchgehend von einer Art: Ein bedientes Ledger ist notwendig
und nicht hinreichend, und eine unbediente Zeile lokalisiert eine Last, statt
eine Behauptung zu widerlegen.


%=====================================================================
\section{Zwei komplementäre Gate-Typen}\label{sec:gates}

Der Vertrags-Audit von Abschnitt~\ref{sec:criterion} ist ein
\emph{logisches} Gate: Er prüft einen festen Text gegen eine feste Liste von
Pflichten, und seine Verdikte sind Aussagen über diesen Text. In dem
Programm, aus dem diese Serie erwuchs, wird er durch ein \emph{empirisches}
Gate ergänzt: ein wegübergreifendes Konsistenzregister (eine \glqq shadow
registry\grqq{} bzw.\ ein Schattenregister), das für jeden untersuchten Weg festhält, welche Aussagen
tatsächlich behauptet werden und wo zwei Wege einander widersprechen würden,
wenn man beide für bare Münze nähme, mit Zulassungsprüfung gegen zertifizierte
Sichten auf einer abgestuften Evidenzskala. Die Instrumentseite dieses
Registers ist im Zookeeper-Papier beschrieben~\cite{zookeeper}; hier ist sie
als komplementäres Gate von Belang und wegen ihrer Grenze, der
\emph{Konsistenzfalle}: Die Übereinstimmung vieler Diagnostiken ist Evidenz
für Kohärenz, niemals für Wahrheit -- ein Register kann konsistent falsch
sein. Die beiden Gates versagen unterschiedlich, und darin liegt ihr
gemeinsamer Wert: Das logische Gate fängt unbediente Pflichten in einem
einzelnen Text; das empirische Gate fängt Widersprüche \emph{zwischen} Texten
und Läufen. Keines ersetzt einen Beweis. Die beiden Gates sind komplementär,
nicht komponierbar: Sie halten getrennte Ausgaben, es gibt keinen gemeinsamen
Entscheidungsalgorithmus, der sie kombiniert, und aus ihnen wird kein
aggregiertes Wahrheitsverdikt und kein Gesamtbestehen gebildet. Eine
Behauptung, die von beiden geprüft wurde, ist auf zwei verschiedene Weisen
geprüft worden und schuldet ihre Mathematik weiterhin.

Abbildung~\ref{fig:dual_gate_architecture} stellt die komplementären Pipelines des logischen und des empirischen Tors gegenüber und veranschaulicht ihre Prüf-Engines, charakteristischen Fehlermodi und die Demarkation gegen eine unzulässige Komposition. Tabelle~\ref{tab:tor_taxonomie} stellt eine synoptische Taxonomie bereit, die das logische Tor ($\mathrm{Br}_\xi$), das empirische Schattenregister, interaktive Theorembeweiser und zertifizierende Algorithmen gegenüberstellt.

\begin{figure}[htbp]
\centering
\resizebox{0.88\textwidth}{!}{%
\begin{tikzpicture}[
  >=Stealth,
  node distance=0.6cm and 0.8cm,
  header/.style={draw=black!80, rounded corners=3pt, line width=0.9pt, align=center, font=\small\bfseries, inner sep=5pt},
  logcol/.style={header, fill=blue!10, draw=blue!70!black, text width=5.6cm},
  empcol/.style={header, fill=teal!10, draw=teal!70!black, text width=5.6cm},
  stepbox/.style={draw=black!60, rounded corners=2pt, line width=0.6pt, align=left, font=\scriptsize, inner sep=4pt, text width=5.6cm},
  alertbox/.style={draw=red!70!black, rounded corners=2pt, line width=0.7pt, fill=red!6, align=left, font=\scriptsize, inner sep=4pt, text width=5.6cm},
  barrier/.style={draw=black!70, rounded corners=3pt, line width=0.8pt, fill=gray!10, align=center, font=\small, inner sep=6pt, text width=12.2cm},
  arrow/.style={->, line width=0.8pt, draw=black!70}
]

  % Headers
  \node[logcol] (h1) at (0, 6.0) {Logisches Tor ($\mathrm{Br}_\xi$)\\[1pt]\footnotesize Einzeltext-Form-Audit};
  \node[empcol] (h2) at (6.8, 6.0) {Empirisches Tor (Schattenregister)\\[1pt]\footnotesize Routenübergreifende Konsistenz};

  % Steps Gate 1
  \node[stepbox, fill=blue!3] (s1_1) [below=0.35cm of h1] {
    \textbf{Gegenstand \& Geltung:}\\[1pt]
    Isolierter Manuskript-Einzeltext $T$
  };
  \node[stepbox, fill=blue!3] (s1_2) [below=0.35cm of s1_1] {
    \textbf{Audit-Instrument:}\\[1pt]
    7 Vertragskomponenten ($C_1$--$C_7$)\\
    BCM-Ledger (v1-T Transport, v1-F Vollinstrument)
  };
  \node[stepbox, fill=blue!3] (s1_3) [below=0.35cm of s1_2] {
    \textbf{Prüf-Engine \& Verdikt-Grammatik:}\\[1pt]
    Hermeneutisches Text-Audit mit 5-Code-Grammatik:\\[1pt]
    $\mathtt{E+}$ (bedient), $\mathtt{E-}$ (unbedient), $\mathtt{I}$ (intern),\\
    $\mathtt{N}$ (extern), $\mathtt{X}$ (nicht-transpositional)
  };
  \node[alertbox] (a1) [below=0.35cm of s1_3] {
    \textbf{Charakteristischer Fehlermodus:}\\[1pt]
    Erkennt keine routenübergreifenden Widersprüche, wenn Einzeltext in sich geschlossen ist.
  };

  % Steps Gate 2
  \node[stepbox, fill=teal!3] (s2_1) [below=0.35cm of h2] {
    \textbf{Gegenstand \& Geltung:}\\[1pt]
    Routenübergreifendes Text- und Rechenkorpus ($A > B > M > C$)
  };
  \node[stepbox, fill=teal!3] (s2_2) [below=0.35cm of s2_1] {
    \textbf{Audit-Instrument:}\\[1pt]
    Wegübergreifendes Konsistenzregister\\
    Automatisierte Behauptungs- und Kollisionsprüfung
  };
  \node[stepbox, fill=teal!3] (s2_3) [below=0.35cm of s2_2] {
    \textbf{Prüf-Engine \& Verdikt-Grammatik:}\\[1pt]
    Empirische Kollisionserkennung:\\[1pt]
    Schattenmatrix, Mehrwege-Kohärenz,\\
    Zulässigkeitsrang auf Evidenzskala
  };
  \node[alertbox] (a2) [below=0.35cm of s2_3] {
    \textbf{Charakteristischer Fehlermodus (Konsistenzfalle):}\\[1pt]
    Gegenseitige Übereinstimmung vieler Wege ist Evidenz für Kohärenz, \textbf{niemals} für Wahrheit.
  };

  % Connectors
  \draw[arrow] (h1) -- (s1_1);
  \draw[arrow] (s1_1) -- (s1_2);
  \draw[arrow] (s1_2) -- (s1_3);
  \draw[arrow] (s1_3) -- (a1);

  \draw[arrow] (h2) -- (s2_1);
  \draw[arrow] (s2_1) -- (s2_2);
  \draw[arrow] (s2_2) -- (s2_3);
  \draw[arrow] (s2_3) -- (a2);

  % Bottom Barrier
  \node[fit=(a1)(a2), inner sep=0pt] (failurepair) {};
  \node[barrier, anchor=north] (bar) at ($(failurepair.south)+(0,-0.45)$) {
    \textbf{Strikte Nicht-Komponierbarkeit \& Demarkation}\\[2pt]
    \footnotesize Die beiden Tore versagen orthogonal $\implies$ getrennter diagnostischer Wert ohne gemeinsamen Entscheidungsalgorithmus.\\[1pt]
    \scriptsize Keines der Tore ersetzt einen mathematischen Beweis; eine von beiden geprüfte Behauptung schuldet weiterhin ihre Herleitung.
  };

  \draw[arrow] (a1.south) -- (a1.south |- bar.north);
  \draw[arrow] (a2.south) -- (a2.south |- bar.north);

\end{tikzpicture}%
}
\caption{Duale Beweisabsicherungs-Architektur: Logisches Tor ($\mathrm{Br}_\xi$) vs.\ Empirisches Tor (Schattenregister).}
\label{fig:dual_gate_architecture}
\end{figure}

\begin{table}[htbp]
\centering
\footnotesize
\renewcommand{\arraystretch}{0.92}
\setlength{\tabcolsep}{3pt}
\caption{Taxonomie von Beweisabsicherungs- und Audit-Toren in der höheren Mathematik.}
\label{tab:tor_taxonomie}\label{tab:gate_taxonomy}
\begin{tabularx}{\textwidth}{@{} >{\raggedright\arraybackslash}p{0.17\textwidth} >{\raggedright\arraybackslash}p{0.19\textwidth} >{\raggedright\arraybackslash}p{0.19\textwidth} >{\raggedright\arraybackslash}p{0.15\textwidth} >{\raggedright\arraybackslash}X@{}}
\toprule
Tor / Paradigma & Audit-Geltung & Prüfer / Trust Base & Ausgabe & Diagnostische Stärke / Grenze \\
\midrule
Logisches Tor ($\mathrm{Br}_\xi$ / BCM) & Einzelner Transportschritt im Text & Hermeneutisch / Fachexperte & 5-Code-Ledger (\texttt{E+}, \texttt{E-}, \texttt{I}, \texttt{N}, \texttt{X}) & Lokalisiert exakte offene Klausel; beurteilt Form, nicht Wahrheit. \\
\addlinespace[4pt]
Empirisches Tor (Schattenregister) & Routen-\allowbreak übergreifende Konsistenz im Korpus & Au\-to\-ma\-ti\-sier\-tes Register ($A>B>M>C$) & Konsistenz- / Kollisionsmatrix & Erkennt Widersprüche zwischen Texten; Konsistenz-Falle. \\
\addlinespace[4pt]
Maschinelle Formalisierung (Lean / Coq) & Deduktiver Kernel über gesamten Beweis & Mechanisierter Proof Assistant & Binäres Gültigkeits-\allowbreak zertifikat & Absolute deduktive Sicherheit; hohe Formalisierungskosten. \\
\addlinespace[4pt]
Zertifizierende Algorithmen (PCC / Laufzeit) & Algorithmisches Eingabe-Ausgabe-Paar & Al\-go\-rith\-mi\-scher Zeugenprüfer & Rechen-\allowbreak technisches Zertifikat & Entscheidbare Instanzprüfung; setzt vorab vereinbarte Policy voraus. \\
\bottomrule
\end{tabularx}
\end{table}

%=====================================================================
\section[Eine protokoll-eingefrorene\texorpdfstring{\newline}{ }retrospektive Zwei-Fall-Illustration (v1-T)]{Eine protokoll-ein\-ge\-fro\-rene re\-tro\-spek\-ti\-ve Zwei-Fall-Il\-lus\-tra\-ti\-on (v1-T)}\label{sec:calibration}

Wir illustrieren das eingefrorene v1-T-Instrument retrospektiv an zwei
historisch entschiedenen Fällen: einem \emph{retrospektiven
Bekannte-Lücke-Fall}, einem Beweisversuch, dessen Lücke später von seinem
Autor bestätigt wurde, und einer \emph{akzeptierten Kontrolle}, einem
akzeptierten Beweis mit einem strukturell vergleichbaren Transportschritt.
Das Protokoll -- Zeilensatz, Fünf-Code-Handbuch, Konfidenzschema,
unabhängige Zweitkodierung mit expliziter Dissensregel,
Rückschaufehler-Kontrollen sowie die Erfolgs- und Falsifikationskriterien --
wurde niedergeschrieben und eingefroren, bevor irgendeine Quelle gelesen
wurde; die Audits haben nichts davon verändert.

Diese zeitliche Priorität ist intern dokumentiert, nicht extern. Das
Protokoll wurde vor der Kodierung nicht öffentlich registriert; seine
zeitliche Priorität ist durch den datierten internen Projektnachweis und die
Auditnotizen dokumentiert. Wir verwenden daher \glqq protocol-frozen\grqq{}
(protokoll-eingefroren) und nicht \glqq extern präregistriert\grqq{}. Wo diese Arbeit
sagt, eine Wahl sei im Voraus festgelegt worden, bedeutet das: prospektiv im
Projektnachweis eingefroren vor den beiden Quellenkodierungsdurchgängen. Eine
heute berechnete Prüfsumme würde die Identität der heute publizierten Datei
feststellen, nicht das Datum, an dem sie eingefroren wurde.

\subsection{Design und Instrumentversion}

Die beiden historischen Fälle wurden mit \textbf{BCM-Ledger v1-T} kodiert,
dem sechszeiligen Transportinstrument, das am 15.~August 2026 vor der
Quellenkodierung eingefroren wurde. Seine versionsspezifischen Zeilen sind
T-R2a, T-R2b, T-R4a, T-R4b, T-R6 und T-R8, mit der in der Tabelle unten
wiedergegebenen fallneutralen Semantik. Dies ist nicht das volle
v1-F-Ledger von Abschnitt~\ref{sec:ledger}; die beiden Versionen sind
ausschließlich durch den publizierten Crosswalk
(Abschnitt~\ref{sec:crosswalk}) miteinander verbunden.

\clearpage
\begin{table}[htbp]
\centering
\small
\caption{Das eingefrorene Sechs-Zeilen-Transportinstrument (BCM-Ledger v1-T) und fallneutrale Prüffragen.}
\label{tab:frozen_instrument_v1t}
\begin{tabularx}{\textwidth}{@{}l >{\raggedright\arraybackslash}X@{}}
\toprule
v1-T-Zeile & Fallneutrale Frage des eingefrorenen Instruments \\
\midrule
T-R2a & Ist das Transport- oder Übertragungsgesetz explizit ausgeschrieben,
        als Formel oder Aussage mit einer Adresse? \\
\addlinespace[3pt]
T-R2b & Werden die transportierten Größen mit einem sichtbaren
        Umrechnungsfaktor mitgeführt? \\
\addlinespace[3pt]
T-R4a & Ist der Kompensations- oder Normierungsmechanismus benannt und
        lokalisiert? \\
\addlinespace[3pt]
T-R4b & Sind seine Existenz und Wohldefiniertheit im Text bewiesen und nicht
        bloß behauptet? \\
\addlinespace[3pt]
T-R6 & Werden zweite Modi oder Fallunterscheidungen dargestellt und
       symmetrisch behandelt? \\
\addlinespace[3pt]
T-R8 & Wird die kettenschließende Kompatibilität hergeleitet statt gesetzt,
       und auf welcher Ebene (lokal oder global summiert)? \\
\bottomrule
\end{tabularx}
\end{table}

Die auditierte Einheit ist in jedem Fall der zentrale Transportschritt,
instanziiert über diese sechs Zeilen. Jede Zeile erhält genau einen der fünf
Codes von Abschnitt~\ref{sec:ledger}. Jeder Fall wurde zweimal kodiert: ein
Erstaudit und eine unabhängige Zweitkodierung, deren Auditor nur das
eingefrorene Protokoll und die Quelle erhielt, nicht die Erstkodierung. Der
Erfolg war im Voraus festgelegt: Im Bekannte-Lücke-Fall muss mindestens eine
Zeile auslösen ($E{-}$/$N$) \emph{und} ihre Adresse muss mit der historisch
dokumentierten Lücke übereinstimmen; in der Kontrolle darf keine Zeile
auslösen (ein $I$ ist unkritisch). Die Fehlschlagsmodi waren symmetrisch
festgelegt: kein Auslösen, Auslösen an der falschen Adresse oder ein Auslösen
in der Kontrolle.

\subsection[Der retrospektive Bekannte-Lücke-Fall:\texorpdfstring{\newline}{ }die Fermat-Ankündigung von 1993]{Der re\-tro\-spek\-ti\-ve Be\-kann\-te-Lü\-cke-Fall: die Fer\-mat-An\-kün\-di\-gung von 1993}

Der erste Fall ist Wiles' Beweisankündigung von 1993, auditiert über den
zeitgenössischen technischen Bericht von Rubin und
Silverberg~\cite{rubinsilverberg}; der Audit betrifft daher den Stand des
Arguments von 1993 \emph{wie berichtet}, eine Geltungsbereichsgrenze, die
ausdrücklich mitgeführt wird. Beide Kodierer erzeugten unabhängig voneinander
dasselbe Bild: Das Transportgesetz und seine Umrechnungsgrößen sind explizit
($\text{T-R2a}, \text{T-R2b} = E{+}$), der kompensierende Mechanismus -- das
Euler-System -- ist benannt und lokalisiert ($\text{T-R4a} = E{+}$), die
Fallunterscheidungen sind dargestellt ($\text{T-R6} = E{+}$), und genau eine
Zeile löst konkordant aus: \emph{die Existenz des benannten Mechanismus}
($\text{T-R4b} = E{-}$), an der Passage, an der der Bericht feststellt, die
Konstruktion eines geeigneten Euler-Systems sei \glqq not yet complete\grqq{}. Die
kettenschließende Zeile T-R8 wurde vom ersten Kodierer mit $E{-}$ kodiert
(als Folge von T-R4b) und vom zweiten mit $I$; nach der Doppelkodierungsregel
unten ist der Dissens verdikt-neutral, da das Erfolgskriterium bereits durch
T-R4b allein erfüllt war.

Die Adresse der auslösenden Zeile ist die Passage, die der Bericht selbst als
unvollständig markiert, und sie ist die Pflicht, die die anschließende
Reparatur des Arguments einlösen musste; die umgebende Kette stand dort nicht
in Frage. Der diagnostische Gehalt liegt ebenso sehr in den vier Zeilen, die
\emph{nicht} ausgelöst haben: Die lokalisierte Lücke ist eng, nicht diffus.
% TODO(operator, P2-7): the sentence above is deliberately anchored only to
% rubinsilverberg, which reports the 1993 state and the announced
% incompleteness but cannot document the 1994--95 repair. If a stronger claim
% about that repair is wanted in the version of record, add a primary or
% contemporaneous anchor for it (candidate to verify: Taylor--Wiles and Wiles,
% Ann. of Math. 141 (1995)) and cite it here. Do NOT strengthen the sentence
% without such an anchor.

\subsection{Die akzeptierte Kontrolle: ein akzeptierter Transportschritt}

Der zweite Fall ist Prasads Volumenformel~\cite{prasad}, deren
Lokal-zu-global-Normierungsschritt eine akzeptierte Instanz derselben
Transportgestalt ist. Beide Kodierer fanden keine auslösende Zeile: fünf
Zeilen in beiden Kodierungen explizit eingelöst; eine Zeile (T-R6) spaltete
sich zwischen $I$ (erster Kodierer, wegen der vom Autor erklärten
Funktionenkörper-Asymmetrie: Die einschlägige Konvergenzaussage wird über
Zahlkörpern durch eine externe Referenz bewiesen, während der Text für globale
Funktionenkörper nur festhält, dass \glqq a similar proof applies\grqq{}
(\cite{prasad}, \S3.2)) und $E{+}$
(zweiter Kodierer), was verdikt-neutral ist, weil $I$ nach dem eingefrorenen
Protokoll unkritisch ist. Der vom Protokoll zugelassene
Geltungsbereichsvorbehalt -- ein Audit eines Transportschrittes innerhalb
eines akzeptierten Beweises statt der ganzen Arbeit -- musste nicht als $X$
aufgerufen werden. Die Kontrolle liefert außerdem ein positives Bild dessen,
was die Zeilen messen: Die Normierungskonstante wird als sichtbarer Faktor
mitgeführt statt stillschweigend gleich eins gesetzt, und die Kette schließt
über eine adressierte Gleichungskette statt über eine Behauptung, sie
\glqq follows formally\grqq{}.
Dies ist die einzige akzeptierte Kontrolle in dieser Arbeit.

\subsection{Die vollständige Kodematrix}\label{sec:matrix}

Beide Kodierungen beider Fälle werden in Tabelle~\ref{tab:code_matrix_retrospective} vollständig berichtet. Die Spalten F1
und F2 sind das Erstaudit und die unabhängige Zweitkodierung des
Bekannte-Lücke-Falls (Wiles 1993, wie berichtet); die Spalten P1 und P2 sind
das Erstaudit und die unabhängige Zweitkodierung der akzeptierten Kontrolle
(Prasad). Keine Zelle unten ist ein Konsenswert.

\begin{table}[htbp]
\centering
\setlength{\tabcolsep}{4.5pt}
\caption{Vollständige doppelkodierte Auditmatrix der retrospektiven Testfälle (Wiles 1993 vs.\ Prasad 1989).}
\label{tab:code_matrix_retrospective}
\begin{tabular}{@{}lcccc@{}}
\toprule
v1-T-Zeile & F1 (Wiles 1.) & F2 (Wiles 2.) & P1 (Prasad 1.) & P2 (Prasad 2.) \\
\midrule
T-R2a & $E{+}$ & $E{+}$ & $E{+}$ & $E{+}$ \\
\addlinespace[3pt]
T-R2b & $E{+}$ & $E{+}$ & $E{+}$ & $E{+}$ \\
\addlinespace[3pt]
T-R4a & $E{+}$ & $E{+}$ & $E{+}$ & $E{+}$ \\
\addlinespace[3pt]
T-R4b & $E{-}$ & $E{-}$ & $E{+}$ & $E{+}$ \\
\addlinespace[3pt]
T-R6 & $E{+}$ & $E{+}$ & $I$ & $E{+}$ \\
\addlinespace[3pt]
T-R8 & $E{-}$ & $I$ & $E{+}$ & $E{+}$ \\
\bottomrule
\end{tabular}
\end{table}

\noindent
Die Übereinstimmung zwischen den Kodierern betrug 10 von 12
Fall-Zeilen-Vergleichen. Die beiden Dissense sind F/T-R8 und P/T-R6; keiner
wurde harmonisiert. Die Regel, unter der diese Kodierungen gelesen wurden,
lautet wie folgt.

\begin{quote}
Ein kritisches Auslösen im Bekannte-Lücke-Fall zählt nur dann für das vorab
festgelegte Erfolgskriterium, wenn beide Kodierer derselben Zeile $E{-}$/$N$
zuweisen. Ein Kontrollfehlschlag wird verzeichnet, wenn einer der beiden
Kodierer $E{-}$/$N$ zuweist. Diskordante Codepaare bleiben sichtbar und
werden niemals durch einen synthetischen Konsenscode ersetzt.
\end{quote}

\noindent
Unter dieser Regel zählt im Bekannte-Lücke-Fall genau eine Zeile als
konkordantes kritisches Auslösen, nämlich T-R4b; in der Kontrolle enthält
überhaupt keine Kodierung $E{-}$ oder $N$. Die Regel ausdrücklich anzugeben
ist eine Offenlegung der bereits angewandten Regel, keine nachträgliche
Änderung der Kodierungen.

\subsection{Verankerung und Verfügbarkeit}\label{sec:anchors}

Jede Kodierung in beiden Audits trägt ein wörtliches Zitat mit einem Seiten-
oder Abschnittsanker aus der benannten Ausgabe, zusammen mit der vom
eingefrorenen Protokoll vorgeschriebenen dreiachsigen Konfidenzangabe:
Quellenbindung, Kriterienpassung und Substanzstatus, jeweils auf einer
dreistufigen Skala. Diese Zeilen-Ankertabellen liegen im datierten
Projektnachweis -- dem eingefrorenen Protokoll, vier Auditnotizen (ein
Erstaudit und eine unabhängige Zweitkodierung je Fall) und der
Operator-Synthese, die den Verdikt-Wortlaut trägt. Dieser Entwurf gibt die
vollständige Kodematrix wieder, nicht aber die Zeilenanker.
% TODO(operator, P1-5): before upload, release the per-row anchor tables (row,
% version/edition, page or section, short quotation, both raw codes, three-axis
% confidence) as a citable supplement and reference it here. Internal file names
% without accessible content do not satisfy this obligation for a public
% version. Note the pipeline rule that _proof-notes/ is private by default: the
% staged release of individual closed notes needs an explicit operator decision.

\subsection{Was die Illustration zeigt -- und was nicht}

Unter den vorab festgelegten Kriterien wurden die beiden vorab festgelegten
Fallkriterien in dieser Illustration erfüllt: Die bestätigte Lücke des
Arguments von 1993 ist an einer einzelnen konkordant auslösenden Zeile
lokalisiert, deren Adresse mit der historisch dokumentierten Lücke
übereinstimmt, und keine Kodierung wies in diesem einen Kontrollfall $E{-}$
oder $N$ zu. Wir formulieren das Ergebnis in der im Protokoll festgelegten
Formulierung: Die beiden Fälle sind \emph{consistent with the certificate's
design intent} (mit der Konstruktionsabsicht des Zertifikats vereinbar). Es
handelt sich um einen retrospektiven Bekannte-Lücke-Fall und eine akzeptierte
Kontrolle -- eine Illustration, keine Validierungsstatistik und keine
Messung irgendeiner Rate. Die Übung ist eine Retrodiktion, keine Vorhersage:
Das historische Ergebnis war beiden Kodierern bekannt, und ihr Wert beruht
vollständig darauf, dass die Zeilenzuordnung vor der Lektüre schriftlich
festgelegt wurde, wobei die falsifizierbaren Alternativen (falsche Adresse,
diffuses Auslösen, ein Auslösen in der Kontrolle) im Voraus angegeben und
nicht eingetreten sind.

Zwei weitere Grenzen gehören hierher. Die Illustration betrifft
ausschließlich v1-T; sie bespielt nicht das volle v1-F-Ledger von
Abschnitt~\ref{sec:ledger}, und nichts weiter unten erbt Stützung von ihr.
Und bei zwei Fällen und zwei Kodierern beschreibt die berichtete
Übereinstimmung dieses Auditpaar; sie ist keine Schätzung der
Zuverlässigkeit des Instruments.

%=====================================================================
\section{Die abc-Instanziierung}\label{sec:abc}

\subsection{Der Gate-Fall}\label{sec:abc-gate}

Dieser Abschnitt berichtet die eine in dieser Arbeit enthaltene
Instanziierung. Es ist eine \emph{selected-rows}-Instanziierung, durchgeführt
mit Zeilen von \textbf{BCM-Ledger v1-F} (Abschnitt~\ref{sec:ledger}): Die
unten bespielten Zeilen sind F-R2a, F-R2b, F-R4a, F-R4b, F-R6a, F-R6b und,
wo ihr Auslöser greift, F-R8; die übrigen Zeilen (F-R0, F-R1, F-R3, F-R5,
F-R7) werden hier nicht kodiert, und eine Kodierung der vollständigen Matrix
für diese Fälle ist auf eine spätere Version zurückgestellt. Die
Zeilenbezeichner unten tragen das Präfix \texttt{F-}, und keiner von ihnen
ist die gleichnummerige Zeile des in Abschnitt~\ref{sec:calibration}
verwendeten eingefrorenen Transport-Ledgers v1-T.

Der strittige Schritt ist ein Teilschritt im Beweis von Korollar~3.12 im
dritten Band der inter-universellen
Teichmüller-Theorie~\cite{Mochizuki2021c}. Seine Schlussfolgerung hat die in
Definition~\ref{def:bridge} axiomatisierte Form: Von einer
log-volumetrischen Ablesung $\nu_q(x)$ des $q$-Pilotobjekts $x \in X_q$ wird
behauptet, sie liege in einer auf der $\Theta$-Seite festgelegten Zielregion
$H_\Theta = \mathbb{R}_{\le -|\log\Theta|}$, wobei die beiden Seiten
konstruktionsbedingt keine gemeinsame Auswertung teilen. Der mit
(xi-e)$\to$(xi-f) bezeichnete Schritt ist die Stelle, an der diese
Zugehörigkeit gezogen wird.

Dass dieser Schritt die Last trägt, ist das Resultat von
Teil~1~\cite{part1} und wird als gegeben genommen. Das Kriterium stellt an
dieselbe Passage eine andere Frage: nicht, wo das Argument dünn ist, sondern
welche Vertragskomponenten die Darstellung bedient. Alles Folgende ist im
Register von Abschnitt~\ref{sec:scope} zu lesen -- eine Zeile, die nicht
liefert, benennt eine Klausel, die ein vollständiges Argument bereitstellen
muss, und sagt nichts darüber, ob sie bereitgestellt werden kann.

\subsection{Der auditierte Teilschritt des ursprünglichen Arguments}\label{sec:abc-original}

Der Audit deckt die am Teilschritt selbst aktivierten Zeilen ab: F-R2a,
F-R2b, F-R4a, F-R4b, F-R6a, F-R6b und, über den Mehrmoden-Auslöser, F-R8.
Die übrigen Zeilen des Ledgers tragen hier kein Verdikt. Die Tabelle unten
sammelt die Verdikte; was folgt, sind die Gründe.

Am auditierten Teilschritt wird keine Transportoperation innerhalb des Links
ausgeführt; F-R2a ist an diesem Teilschritt daher X, nicht E+. Die von der
Schlussfolgerung benötigte rahmenwerkübergreifende numerische
Identifikation wird gesondert unter F-R2b und F-R6a beurteilt. Eine
gemeinsame Auswertung in C löst Komponente~6 nicht ein, solange nicht die
verlangten Kompatibilitätsidentitäten vorhanden und adressiert sind. Was die
erste dieser Aussagen stützt, ist, dass der Paralleltransport-Mechanismus
als einer beschrieben wird, der \glqq does not consist of a simple instance of
transport of some set-theoretic region\grqq{}, und dass die Ablesung auf der
$q$-Seite als \glqq coric\grqq{}/\glqq fixed\grqq{} erklärt und ausdrücklich nicht den
Unbestimmtheiten unterworfen wird, die die $\Theta$-Seite trägt: Die Ablesung
behält ihre Identität durch Festsetzung und durch das Fehlen eines
Transportschrittes, nicht durch eine bewiesene Relation.

F-R2b liefert teilweise, und nicht mehr als das. Bei (xi-d) werden die
beiden Seiten in eine Auswertung gebracht, wodurch erst entsteht, was der
Text \glqq completely comparable objects\grqq{} nennt, und globale arithmetische Grade
werden auf beiden Seiten als Log-Volumina gelesen. Das ist eine
Identifikation, die von einer geteilten Lesekonvention getragen wird und
nicht von einem sichtbaren Umrechnungsfaktor oder einer
Kompatibilitätsgleichung, und der Crosswalk deckelt eine Zeile dieser Gestalt
bei \texttt{I}. Die Zeile erreicht diese Obergrenze und fällt nicht darunter,
weil die Identifikation an einem benannten Teilschritt selbst ausgewiesen ist
und am Text abgelesen werden kann: Was der dokumentierten Rekonstruktion
fehlt, ist eine Herleitung der Identifikation, nicht die Identifikation. An
diesem Schritt ist die Zeile überdies nicht von F-R6a trennbar, da die
Identifikation der beiden Auswertungscontainer Teil des Brückenschrittes
selbst ist.

F-R6a liefert nicht. Eine numerische Auswertung $\nu_C$ auf dem Container
wird ausgewiesen; was die gehärtete Komponente~6 zusätzlich verlangt -- die
von der Behauptung benötigten Kompatibilitätsidentitäten, insbesondere
$\nu_C \circ i_q = \nu_q$ auf dem angegebenen Quellbereich --, wird am
auditierten Schritt nicht hergeleitet. Die quellseitige Ablesung wird durch
Festsetzung fixiert gehalten -- die oben unter F-R2a zitierte Passage, die
die Ablesung als \glqq coric\grqq{}/\glqq fixed\grqq{} und als von den Unbestimmtheiten der
$\Theta$-Seite ausgenommen erklärt --, und eine festgesetzte Invarianz ist
nicht die Identität, nach der Komponente~6 verlangt. Die Zeile würde nur dann
teilweise liefern, wenn eine dokumentierte Rekonstruktion dieser Identität
bereitgestellt würde; dieser Audit stellt keine bereit. Die frühere Lesart
dieser Passage, die Komponente~6 als bedient verzeichnete, beruhte auf der
schwächeren Entweder-oder-Formulierung, die Definition~\ref{def:bridge} nicht
mehr enthält.

F-R6b liefert nicht, und sie ist die wichtigste kettenschließende Pflicht.
Die Zugehörigkeit wird nicht aus der zulässigen Verknüpfung hergeleitet,
sondern aus einer den zulässigen Operationen auferlegten Bedingung gewonnen:
Die Inklusion \glqq then follows formally\grqq{} aus der Forderung, dass die
Konstruktion der Ausgabemöglichkeiten insbesondere eine Konstruktion der
Eingabe-Ablesung darstellt, und eine Zusammenfassung desselben Arguments
spricht von \glqq two tautologically equivalent ways\grqq{}, ein Log-Volumen zu
berechnen. Wo das benannte Ordnungslemma von Komponente~7 stehen würde -- mit
seinen Hypothesen, einer Behauptung der Anwendbarkeit an diesem Schritt und
einer Beweisadresse --, steht stattdessen ein Paar erklärter Eigenschaften
der Algorithmusausgabe. Die Zeile verzeichnet das Fehlen eines solchen
adressierten Lemmas. Sie verzeichnet kein Urteil darüber, ob eines bewiesen
werden könnte.

Zwei Fairness-Aussagen gehören zu diesem Verdikt und nicht daneben. Die
Konfidenz ist gerade deshalb hoch, weil die tragenden Passagen diese Wörter
selbst verwenden; und \glqq tautological\grqq{} wird dort affirmativ gebraucht, wobei
der Text an anderer Stelle seine Schlussfolgerungen als nichttrivial, wenn
auch im Wesentlichen tautologisch bezeichnet. Für diesen Autor ist die
Tautologizität der Punkt, kein Eingeständnis, und sie als Selbstwiderlegung
zu lesen wäre ein umgekehrtes Autoritätsargument.

Die Lastzeilen verlangen die zweite Aufspaltung. Im auditierten Band ist die
von den Theta-Werten $q^{j^2}$ getragene Last benannt, aber nicht
quantifiziert -- was an der identifizierten Stelle gebucht wird, sind
Etiketten und nicht Größenordnung --, und am Brückenschritt wird die interne
Struktur bewusst beiseitegelassen. Die Quantifizierung wird nicht
ausgelassen, sondern angekündigt, ihre Adresse ist als die explizite
Log-Volumen-Rechnung in \S1 des vierten Bandes~\cite{Mochizuki2021d}
angegeben, und diese Adresse zahlt sich aus: Die Last wird durch jene
Rechnung bis zur geschlossenen Form $(\ell+1)/24$ getragen und steht in der
finalen Schranke als
$-\tfrac{1}{6}\bigl(1-12/\ell^{2}\bigr)\log q$. F-R4a liefert daher für den
Verbund der beiden Bände; F-R4b nicht, da jener Abschnitt die
Schlussfolgerung des Korollars als Eingabe nimmt und nachgelagert zur Brücke
quantifiziert. Beide Zeilen werden an Komponente~5 berichtet, wo der Crosswalk
das Lastbudget lokalisiert. Eine frühe Sorge, festgehalten in der
Ledger-Notiz vom Mai 2026, die $j^2$-Faktoren könnten in einer Koordinate
verschwinden und ungebucht zurückkehren, hält für den Verbund nicht stand und
wird zurückgezogen: anderswo wie angekündigt gebucht ist nicht derselbe
Befund wie ungebucht.

F-R8 wird hier ausgelöst, da das Bereichsmodell der Behauptung einen
Alternativmodus enthält, der die Ordnungsrelation betrifft. Sie liefert
teilweise: Der zweite Modus ist benannt, der Vertikalverschiebungskonflikt
angegeben und sein Preis erklärt, aber er wird vorab in eine einseitige
Ungleichung überführt und am Schritt selbst suspendiert.

Die auditierten Zeilen reduzieren sich daher nicht auf eine einzelne übrig
bleibende Pflicht. F-R2a ist am Teilschritt außerhalb des Geltungsbereichs,
F-R2b liefert teilweise, und F-R4a liefert für den Verbund der beiden Bände.
Gegen dieses Profil werden die verbleibenden Pflichten wie folgt gewichtet.
Der Audit unterscheidet eine wichtigste kettenschließende Pflicht von drei
weiteren Befunden: F-R6b ist die wichtigste Ordnungslemma-Pflicht; F-R6a
verzeichnet die fehlende Kompatibilitätsidentität der gemeinsamen
Auswertung; F-R4b verzeichnet ein gesondertes Defizit der Buchung an der
Brücke; und das ausgelöste F-R8 verzeichnet eine teilweise Behandlung des
zweiten Modus. Das Ledger lokalisiert daher unterschiedlich gewichtete
Befunde statt einer einzelnen übrig bleibenden Zeile.

\subsection{Ein vergleichender Diskriminierungsfall: Joshis
ATS-Programm}\label{sec:abc-control}

Ein Kriterium, das überall dasselbe Verdikt zurückgibt, ist eine Schablone,
kein Instrument. Wir wenden daher dieselben ausgewählten v1-F-Zeilen, mit
unveränderter Zeilensemantik, auf ein zweites Programm desselben
\emph{dispute complex} an, das dieselbe Ungleichung auf anderem Wege
erreicht: Joshis arithmetische
Teichmüller-Räume~\cite{JoshiATSIII,JoshiATSIV}. Dieser Fall ist keine
Positivkontrolle und liefert keine unabhängige Bestätigung der Ungleichung.
Er ist ein methodisch andersartiger, vorläufiger Vergleich, mit dem gefragt
wird, ob dieselben ausgewählten Zeilen ein anderes Muster bedienter,
rekonstruierter und unbedienter Pflichten erzeugen. Die einzige akzeptierte
Kontrolle in dieser Arbeit ist die in
Abschnitt~\ref{sec:calibration}. Sie tun es: Das Profil unterscheidet sich,
und die verbleibende Pflicht hat eine andere Adresse und eine andere Gestalt.

F-R2a ruht auf einer bewiesenen Relation statt auf einer Festsetzung: Die
Skalierungsrelation zwischen den Bewertungen entlang des programmeigenen
Links ist dort ein Satz, und der $j^2$-Faktor wird sichtbar in eine explizit
definierte log-volumetrische Skala getragen. Hier trennt die Aufspaltung
sauber, weshalb dieser Audit sie hervorgebracht hat. F-R2b hingegen ruht auf
der Wahl derselben globalen Normierung auf beiden Seiten -- einer
Konvention, nicht einem ausgewiesenen Umrechnungsfaktor -- und liefert daher
höchstens teilweise, genau so, wie der Crosswalk eine Zeile dieser Gestalt
deckelt. Sie erreicht diese Obergrenze, weil die Normierung ausdrücklich
gewählt und angegeben ist, sodass die Rekonstruktion dokumentiert ist; was
fehlt, ist ein sichtbar durch den Schritt getragener Umrechnungsfaktor.

Die Lastpflicht wurde auditiert, bevor die Aufspaltung eingeführt wurde, und
wird daher als ein einziges Verdikt berichtet, bei F-R4a. Die Last wird gegen
eine globale Nebenbedingung gebucht, die Produktformel, mit hoher Konfidenz
für die Adresse und mittlerer für die Vollständigkeit, wobei die
kompensierende Renormierung erschlossen und nicht konstruiert ist. Die
Adresse ist bedient und die Vollständigkeit ist nicht etabliert, sodass die
Zeile für die Adresse liefert und für die Vollständigkeit nur teilweise --
nicht uneingeschränkt. Ein Audit an der Aufspaltung müsste die bediente
Adresse von der erschlossenen Existenz trennen; dieser hat das nicht getan,
und F-R4b trägt daher in diesem Fall keinen gesonderten Befund.

F-R8 ist so geliefert, wie die Zeile es verlangt: Eine globale
Frobenius-Verschiebung ist benannt, für die beiden Schranken asymmetrisch
eingesetzt und vom Autor als Korrektur seiner eigenen früheren Version
gekennzeichnet -- und weil sie sichtbar ist, hat die verbleibende Pflicht
überhaupt eine Adresse.

F-R6a liefert teilweise, und die Gestalt des Rests ist der Ertrag dieses
Falls. Die Ordnungsrelation, die die Schlussfolgerung trägt, gilt zwischen
Summen über alle Stellen; stellenweise wird sie ausdrücklich zurückgewiesen,
wobei der Autor angibt, dass lokale termweise Vergleiche an jeder Primstelle
nicht angestellt werden können, ohne die gleichzeitig geltenden Normierungen
zu verletzen, und die Bemerkung des Satzes selbst festhält, dass diese Wahl
\glqq precludes direct comparison of various local quantities\grqq{}. Die
Identifikation der beiden Auswertungen ruht auf identischer globaler
Normierung und wird an dieser Stelle eine Tautologie genannt. Dieses Wort
muss in seinem Geltungsbereich präzisiert werden: Es kommt an mehreren Stellen des Programms mit
unterschiedlicher Reichweite vor, nur seine Verwendung für die mittlere
Gleichheit des tragenden Satzes ist hier tragend, und es ist nicht so, dass
der Autor seinen eigenen arithmetischen Beweis des Korollars eine Tautologie
nennt.

Bei F-R6b unterscheidet sich dieser Fall der Form nach am schärfsten vom
ursprünglichen Argument, und hier ist die Grenze des Auditors am
sorgfältigsten einzuhalten. Was die Zeile prüft, ist vorhanden: ein benannter
Satz mit angegebenen Hypothesen, eine Behauptung seiner Anwendbarkeit am
Schritt und ein ausgeschriebener Beweis mit einer Adresse. Das ist alles, was
ein bedientes F-R6b zertifiziert. Dem steht ein Substanzbefund gegenüber,
hier bei F-R6b berichtet, weil die Pflicht, die er betrifft, dort lokalisiert
ist: Der Substanz-Audit von Teil~5~\cite{part5} typisierte den Kernschritt
jenes Beweises gegen die von ihm zitierte Konstruktion als offen, sodass der
Kern der Pflicht offen bleibt. Dieser Befund ist kein Formverdikt, und der
Auditor entscheidet ihn nicht.

Zwei Korrekturen revidierten die historische Lesart dieses Falls in
entgegengesetzte Richtungen. Innerhalb des bestimmten Normierungsrahmens,
der im auditierten Teil~4~\cite{part4} benannt ist, wurde die Identifikation
als definitorisch statt als verdeckte Wahl gelesen; dies ist keine
unbedingte Identifikation über ATS-Felder, Stellenbereiche oder
Volumenkonstruktionen hinweg. Mengenstabilität und die Festlegung des
Maßes bleiben gesonderte Pflichten. Der Teil-5-Substanzbefund~\cite{part5}
ist von dieser Formlesart getrennt. Das im historischen Zusammenhang
zitierte Starrheitsresultat~\cite{part6} gilt unter seinen angegebenen
Feld-, Allstellen- und Normierungshypothesen; es beseitigt keine dieser
Hypothesen und liefert keinen unbedingten ATS- oder Volumenschluss.
Tabelle~\ref{tab:abc_ats_comparison} erhält die historischen vergleichenden
Auditergebnisse, begrenzt auf die auditierten Versionen. Der Band, der die
abc-Schlussfolgerung trägt, bezeichnet sich als vorläufige Version für
Kommentare.

\begin{table}[htbp]
\centering
\small
\caption{Vergleichendes Auditprofil für den abc-Transportschritt (ursprüngliche IUT-Darstellung vs.\ ATS-Programm).}
\label{tab:abc_ats_comparison}
\begin{tabularx}{\textwidth}{@{}l >{\raggedright\arraybackslash}X >{\raggedright\arraybackslash}X@{}}
\toprule
v1-F-Zeile & Ursprüngliches Argument & Vergleichsfall (ATS-Programm) \\
\midrule
F-R2a & außerhalb des Geltungsbereichs (X): keine Transportoperation
     innerhalb des Links am Teilschritt
   & liefert (hoch): eine bewiesene Skalierungsrelation entlang des
     programmeigenen Links \\
F-R2b & liefert teilweise: Identifikation durch eine geteilte
     Lesekonvention; hier nicht von F-R6a trennbar
   & liefert teilweise: dieselbe globale Normierung, durch Konvention
     gesetzt \\
F-R4a & liefert (hoch) für den Verbund der beiden Bände
   & unaufgespalten auditiert, hier berichtet: liefert für die Adresse,
     teilweise für die Vollständigkeit -- die kompensierende Renormierung
     ist erschlossen und nicht konstruiert \\
F-R4b & liefert nicht (hoch): nachgelagert zur Brücke quantifiziert
   & kein gesonderter Befund: nicht an der Aufspaltung auditiert (siehe
     F-R4a) \\
F-R6a & liefert nicht: keine Kompatibilitätsidentität am Schritt hergeleitet
   & liefert teilweise (hoch): Ordnung nach Summation, stellenweise
     zurückgewiesen \\
F-R6b & liefert nicht (hoch): Zugehörigkeit \glqq follows formally\grqq{}
   & textlich bedient: benannter Satz, Hypothesen, Beweisadresse;
     Kernschritt von einem Substanz-Audit als offen typisiert \\
F-R8 & liefert teilweise (hoch): benannt und bepreist, am Schritt suspendiert
   & liefert (hoch): benannt, asymmetrisch, erklärt \\
\bottomrule
\end{tabularx}
\end{table}

\noindent
Die auditierte Einheit waren in beiden Fällen diese Zeilen am identifizierten
Schritt; die übrigen Zeilen von Abschnitt~\ref{sec:ledger} tragen hier kein
Verdikt, und eine mit X kodierte Zeile wird keiner der beiden Darstellungen
angelastet. Was der Vergleich lokalisiert, ist ein Unterschied in der
\emph{Form} der verbleibenden Pflicht: ein Paar erklärter Eigenschaften,
verteilt über eine Bemerkung und mehrere Teilschritte, gegen eine einzelne
benannte Gleichheit mit ausgeschriebenem Beweis. Ob dieser Beweis korrekt
ist, ist keine Frage, die dieses Instrument beantwortet.

\subsection{Vererbung: was mit einer Zitation mitreist}\label{sec:abc-inheritance}

Ein nachgelagertes Vererbungsbeispiel führt die Frage über den auditierten
Schritt hinaus. Zwei Anwendungsarbeiten von
Zhou~\cite{ZhouI,ZhouII} leiten diophantische Konsequenzen her, die auf der
Ungleichung von Korollar~3.12 ruhen, und der Audit stellte eine Frage: wie
die Brücke übernommen wird. In der ersten erscheint sie als Proposition,
deren gesamter Beweis ein Satz ist -- \glqq The proposition follows from the
$\mu_6$-version of [37], Corollary~3.12\grqq{}, wobei sich die Referenz auf den
dritten Band in der Nummerierung jener Arbeit bezieht -- und wird genau
einmal verwendet, wonach die nachgelagerte Kette weiterläuft. In der zweiten
wird das Korollar in der auditierten Version überhaupt nicht benannt: Die
Brücke wird über die erste Arbeit vererbt und auf eine neue Klasse von
Anfangsdaten ausgedehnt, indem eine einzelne Hypothese durch die zitierten
Bände verfolgt wird, ohne dass der Brückenschritt berührt würde. Was beide
Arbeiten an Eigenem beitragen, liegt jenseits der Brücke.

Nichts davon ist ungewöhnlich: Ein etabliertes Resultat zu zitieren, statt es
neu herzuleiten, ist genau der Zweck der Zitation. Was den strukturellen
Punkt sichtbar macht, ist, dass beide Arbeiten klar angeben, wo sie einen
eigenen Schritt auslassen. Diese Vorbehalte haften an ihren eigenen
Auslassungen, und keiner haftet am importierten Schritt -- nicht, weil die
Autoren es abgelehnt hätten, einen zu machen, sondern weil an einer Zitation
nichts einen Vertragszustand trägt. Eine bediente Zeile und eine unbediente
Zeile reisen identisch durch eine Bibliographie.

Das ist das praktische Argument für ein Vertragszertifikat als
\emph{Zitationsschnittstelle}. Ein auditiertes Ledger ist ein zitierbares
Artefakt: Es gibt Zeile für Zeile an, was das zitierte Resultat als bedient
und was es als offen gelassen befunden wurde, sodass eine importierende
Arbeit einen bekannten statt eines unbekannten Vertragszustands erbt. Unter
einer solchen Disziplin -- in Teil~3~\cite{part3} als Import-Ledger
formuliert -- würde sich an der Mathematik keiner der beiden Arbeiten etwas
ändern; jede trüge eine weitere Zeile, die die Zeile benennt, die ihr Import
unbedient lässt. Beide Befunde sind im Geltungsbereich auf die auditierten Versionen beschränkt,
wobei die zweite Arbeit sich selbst als ältere, zu aktualisierende Version
bezeichnet, und eine Frage, die der Audit offen ließ -- ob ein
Normierungsfaktor in der Aussage der ersten Arbeit dieselbe Größe
renormiert, die das Korollar trägt --, bleibt auch hier offen.

\subsection{Ein externer Vergleichspunkt}\label{sec:abc-external}

Der letzte Punkt dieses Abschnitts stammt nicht von uns. Ein Zwischenbericht
eines Formalisierungsprojekts, publiziert im Juli 2026~\cite{LANA2026},
isoliert an derselben Passage ein Kompatibilitätsproblem -- und isoliert es
in der Gestalt eines Vertrags. Seine dort mit (9-1) nummerierte Bedingung
verlangt die Existenz einer geeigneten Wahl der ganzzahligen Struktur, unter
der zwei Auswertungsisomorphismen mit gemeinsamem Ziel übereinstimmen: der
eine entstehend aus dem $q$-Piloten in seiner nativen Struktur, der andere
anabelsch und kummertheoretisch gewonnen und Unbestimmtheiten unterworfen.
Der Bericht stellt fest, dass das Projekt \glqq has not yet reconstructed a proof
of the required compatibility\grqq{}, und dass der Beweis der finalen numerischen
Ungleichung \glqq hinges on the compatibility (9-1) which is not manifestly
false. However, we, the LANA project, do not have a proof of (9-1) at this
time.\grqq{}

Gegen Definition~\ref{def:bridge} gelesen ist die vorgeschlagene Zuordnung
objektweise, kein verifizierter Transfer: Ein gemeinsamer Zielquotient,
typisierte Abbildungen, Ablesungen, Ganzheitsstrukturen und Unbestimmtheiten
müssten mit ihren tatsächlichen Objekten, Domänen, Quantoren und ihrer
Implikationsrichtung identifiziert werden. Die Existenz eines Quotienten oder
eines geteilten Ziels allein bedient nicht die Kompatibilitätsidentitäten
von Komponente~6. Ist eine Ablesegleichheit im identifizierten
gemeinsamen Werteraum $V$ bewiesen, liefert Reflexivität den numerischen
Vergleich in $V$; nicht umgekehrt. Dies liefert für sich nicht die
Containerordnung (H2): $\nu_C$ muss Ordnung nicht reflektieren. Gleichheit
der identifizierten Bilder in $C$ liefert (H2) durch Reflexivität in $C$;
numerischer Vergleich allein verlangt Ordnungsreflexion oder einen gesondert
adressierten Beweis von (H2). Zugehörigkeit zu
$H_\Theta$ verlangt zusätzlich die anwendbaren (H1)--(H5), einschließlich
der Zielschranke (H5). Dieser Vergleich beweist keine solche IUT-zu-C7-Brücke.

Die Ledger-Notiz, die das Kriterium im Mai 2026 fixierte, hielt eine bedingte
Erwartung fest: Sollte eine Formalisierungsanstrengung diese Passage tragen,
müsste der Vertrag in ihr als Lemma auftreten. Der Bericht hat diese Gestalt.
In diesem einzelnen Zwischenbericht, an dieser einzelnen Passage, besetzt
eine unabhängig formulierte Kompatibilitätsbedingung eine zu F-R6b analoge
Position. Wir verzeichnen eine nachträglich festgestellte strukturelle
Ähnlichkeit, keine Konvergenz auf das allgemeine Kriterium und keine
Bestätigung unseres Fallverdikts. Ein Prioritätsanspruch wird ebenfalls nicht
erhoben; die Analyse ist die jenes Projekts und in seinen eigenen Begriffen.
Die Grenzen sind ebenso deutlich: Das Dokument ist ein Zwischenbericht ohne
abschließendes Verdikt und hält internen Dissens fest, und unsere Zuordnung
ist objektweise, kein Beweis, dass (9-1) und Komponente~7 äquivalent sind. Die
eigene Position des Berichts -- verlangt, unbewiesen, nicht offensichtlich
falsch -- ist genau das, wofür die Verdikt-Grammatik von
Abschnitt~\ref{sec:ledger} existiert: Eine Zeile, die nicht liefert,
lokalisiert eine Last und widerlegt nichts.

%=====================================================================
\section{Ausblick: RH-Uniformitätsbehauptungen}\label{sec:rh}

Das Kriterium ist nicht spezifisch für die abc-Literatur.
Positivitätszugänge zur Riemannschen Vermutung enthalten dieselben zwei
Gestalten -- Transportschritte (Resultate endlicher Modelle, zu einem
Grenzobjekt getragen) und global-only-Behauptungen (Positivität erst nach
Summation verfügbar) --, und die entsprechenden Ledger-Zeilen könnten dort
instanziiert werden: Uniformitätspflichten für Schranken endlicher
Abschnitte, Quotientenpflichten dort, wo Majorisierungsargumente
ausgeschlossen sind, und eine explizite Zahlungspflicht für archimedische
Masse. Wir stellen jede solche Instanziierung auf eine eigene Behandlung
zurück, synchronisiert mit dem Selbstanwendungsteil dieser Serie, weil dort
unser eigenes Programm zu den auditierten Objekten gehört und die
Protokoll-Einfrierungsdisziplin von Abschnitt~\ref{sec:calibration} verlangt,
dass die operationalisierten Zeilen, Primärquellen und Behauptungsgrenzen
festgelegt sind, bevor der Abschnitt geschrieben wird. Dieser Abschnitt ist
ein Ausblick und keine Instanziierung: Die vorliegende Arbeit enthält bewusst
keine RH-seitigen Verdikte.

\section[Eine verwandte Struktur: Global-only-Behauptungen]{Eine verwandte Struktur:\texorpdfstring{\\}{ }Global-only-Behauptungen}\label{sec:global}

Eine wiederkehrende Gestalt in beiden Instanziierungsbereichen ist die
\emph{global-only}-Behauptung: eine Aussage, von der behauptet wird, sie
gelte erst nach Summation über alle Stellen, wobei lokaler, stellenweiser
Zugriff ausdrücklich nicht verfügbar ist. In Weil-Positivitätszugängen zur RH
ist die Positivität einer quadratischen Form eine globale Summe, deren
einzelne Terme kein Vorzeichen tragen; in der abc-Literatur beschränkt die
Quelle des wichtigsten Reparaturprogramms ihre Vergleiche mit ausdrücklichen
Worten auf die summierte Ebene (\glqq one cannot make local term by term
comparisons at each prime\grqq{}~\cite{JoshiATSIII,JoshiATSIV}), und ihr
Normierungsmechanismus ist eine Produktformel-Nebenbedingung -- wiederum
eine Aussage über ein Produkt über alle Stellen, nicht über irgendeinen
einzelnen Faktor. Die Auditspur für beide Lesarten findet sich in den
Teilen 4 und 5 dieser Serie~\cite{part4,part5}.
% TODO(operator, P2-7): the quotation above is now anchored to the primary ATS
% source rather than to our own audit parts, but it still lacks a precise
% locator. Add the volume, section or page of ATS III/IV for this sentence
% before upload; the same locator is also wanted at the parallel passage in
% Section \ref{sec:abc-control}.

Für das Kriterium löst diese Gestalt eine spezifische Pflicht aus. Wo eine
Behauptung nur auf der summierten Ebene lebt, können die
Vertragskomponenten 6--7 nicht stellenweise bedient werden; was dann
bereitgestellt werden muss, ist ein \emph{Vertragslemma}. Wir geben an, was
ein solches Lemma enthalten muss.

\begin{definition}[Global-only-Vertragslemma]\label{def:globallemma}
Es sei $(x_v)_{v\in\Sigma}$ eine zulässige Familie, $G((x_v)_v)=e$ eine
angegebene globale Nebenbedingung und $S((x_v)_v)$ das von der Behauptung
verwendete Aggregat. Ein Global-only-Vertragslemma ist eine benannte,
adressierte Aussage, die (G1) die zulässige Klasse und die Konvergenz- bzw.\
Bereichsbedingungen, (G2) die exakte globale Nebenbedingung, (G3) die für die
Behauptung benötigte Aggregatgleichheit oder Ordnungsimplikation und (G4) die
lokalen Freiheitsgrade angibt, die nach Auferlegung der Nebenbedingung
verbleiben. Es löst F-R6b nur auf globaler Ebene ein. Eine stellenweise
Starrheitsschlussfolgerung verlangt eine zusätzliche benannte Injektivitäts-,
Striktheits- oder Eindeutigkeitshypothese und darf nicht aus der Globalität
allein erschlossen werden.
\end{definition}

In dieser Arbeit ist die Definition eine Ausblicksschnittstelle, keine zweite
Instanziierung; es wird kein global-only-Ledger-Verdikt berichtet. Das
Werkzeug-Supplement dieser Serie~\cite{gnc} führt eine Konstruktion dieser
Art für die globale Normierung des abc-Reparaturprogramms aus, und der
Starrheitssatz von Teil 6~\cite{part6} weist einen Fall auf, in dem eine
globale multiplikative Nebenbedingung über alle Stellen zusammen mit den dort
angegebenen Hypothesen sehr wohl stellenweise Starrheit erzwingt -- was eine
Instanz von unter benannten Zusatzhypothesen eingelöstem (G4) ist und keine
allgemeine Lizenz, Starrheit an der Globalität abzulesen. Innerhalb von v1-F
wird die ausgelöste Zeile F-R8 auf globaler Ebene nur dann aktiviert, wenn
mehrere globale Modi oder Zweige die Ordnung des Aggregats beeinflussen
können. Die Lehre für die Auditierung ist symmetrisch: global-only ist kein
Mangel, aber es verwandelt lokale Pflichten in eine einzelne benannte
globale -- und der Audit muss prüfen, dass diese eingelöst und nicht in
Notation aufgelöst wird.

\subsection{Aktuelle Entwicklungen in Teilen 4--6}\label{sec:current-series}

Die späteren v0.3-Papers sind zusätzliche versionsspezifische
Forschungsreferenzen, keine Ersatzfassungen der v0.2-Bindungen des
vorstehenden retrospektiven Audits. Teil~4~\cite{part4current} erhält eine
teilweise dokumentierte Normierungskonvention und trennt die typisierten
Logbild-, Hüllen- und Logvolumenverträge. Teil~5~\cite{part5current}
unterscheidet rohe und normierte Koordinaten und verlangt für seine
Frobeniusroute die tatsächlichen gemeinsamen fixierten Koordinaten,
Koeffiziententransformation und Periodenbildaktion. Er dokumentiert auch
die tatsächlich gedruckten Betragsstriche im eng geprüften ATS-III-Korollar;
die frühere atomare Vorzeichenentlastung wird unter den dort ausdrücklich
genannten Tate-Stellen- und Endlichkeitsannahmen zurückgenommen. Diese
späteren Befunde kodieren die historischen Zeilen nicht um und entscheiden
nicht die gesamte ATS-, abc- oder IUT-Kette.

In der Standardmodul-Normierung von Teil~6~\cite{part6current} sind die
lokalen Gradfaktoren enthalten, sodass
\[
 \sum_{v\in M_L}\log b_v(x)=0\qquad(x\in L^\times).
\]
Für ein festes Zahlfeld $L$, alle seine Stellen einschließlich der
archimedischen, reelle von $x$ unabhängige Gewichte und jedes
$x\in L^\times$, mit lokaler Normrealisierung $\rho_v=b_v^{a_v}$,
betrifft die gewichtete Produktformel die endgültigen
Effektivgewichte
\[
 \beta_v=a_v\lambda_v,\qquad
 \sum_{v\in M_L}\beta_v\log b_v(x)=0\quad\text{für alle }x\in L^\times.
\]
Unter diesen Hypothesen liefert Starrheit $\beta_v=c$; sie identifiziert
die rohen $\lambda_v$ nicht mit $c$, wenn Gradfaktoren getrennt stehen.
Ein konstanter globaler Faktor, einschließlich $j^2$, ist zulässig.
Auf $L^\times$ beschränkte Tests an Stellen einer Erweiterung $L'$
bestimmen nicht die einzelnen $L'$-Gewichte. Gemeinsame Feld-, Norm- und
Freeze-Identifikationen bleiben Voraussetzungen eines Zweigvergleichs.

Für einen Größenverbraucher wird eine lineare Familie verlangt:
\[
 Q_y=\sum_v\beta_{y,v}q_v,\qquad Q_{\rm ref}=\sum_v q_v,
 \qquad Q_y=c_yQ_{\rm ref},
\]
mit von $y$ unabhängigen $q_v$ und endlichen Summen oder gerechtfertigter
absoluter Konvergenz. Die letzte Gleichheit gilt bedingt unter denselben
Effektivgewichtshypothesen; $Q_{\rm ref}=Q_{\rm cl}$ ist eine zusätzliche
klassische Referenzbrücke. Punktweise Normierung (5.3.3) liefert unter ihrer
angegebenen Identifikation den klassischen Grad auf $L^\times$ und die
Höhe auf $\mathbb P^n(L)$, ohne automatische Übertragung auf algebraische
Erweiterungen oder Geradenbündelhöhen. Mitgliedschaft und Nichtmitgliedschaft
von Maßen und Logvolumina in dieser Klasse bleiben unbekannt; eine beliebige
ATS-Hülle wird nicht stillschweigend mit einer Maximalordnungs-Modulhülle
identifiziert. Dies sind aktuelle Geltungsbereichspräzisierungen, keine
neuen F-R6-Verdiktscodes.

%=====================================================================
\section{Grenzen und Fairness}\label{sec:limits}

Vier Grenzen umschließen alles Vorstehende. \emph{Erstens} beurteilt das
Kriterium Form. Kein Verdikt in dieser Arbeit oder in der Serie ist ein
Wahrheitsverdikt: nicht über abc, nicht über IUT, nicht über irgendein
auditiertes Programm. Eine unbediente Ledger-Zeile ist eine Pflicht mit einer
Adresse; Theoreme werden durch Mathematik widerlegt, nicht durch Audits. Was
der Auditor zertifiziert, ist noch enger: dass eine Pflicht in einer benannten
Version vorhanden, typisiert und adressierbar ist, niemals, dass die
Mathematik, die sie benennt, tragfähig ist. \emph{Zweitens} ist der
Auditrahmen in der zugrunde liegenden Kontroverse neutral. Dieselben Zeilen
werden auf jeden auditierten Text angewandt -- das ursprüngliche Programm,
die Lesart seiner Kritiker und das Reparaturprogramm --, und die
retrospektive Illustration von Abschnitt~\ref{sec:calibration} wurde gerade
deshalb mit einem symmetrischen Falsifikationskriterium entworfen, damit das
Instrument hätte scheitern können. \emph{Drittens} ist die Evidenzbasis
konstruktionsbedingt klein und versionsspezifisch abgegrenzt: ein retrospektiver
Bekannte-Lücke-Fall und eine akzeptierte Kontrolle unter v1-T, eine
abc-Instanziierung mit einem vergleichenden Diskriminierungsfall unter v1-F,
dazu ein nachgelagertes Vererbungsbeispiel und ein externer Vergleich. Nichts
in der v1-F-Instanziierung erbt Stützung von der v1-T-Illustration. Wir
verwenden die im Protokoll festgelegte Formulierung \glqq consistent with design
intent\grqq{} (mit der Konstruktionsabsicht vereinbar) und vermeiden durchgehend
das Wort \glqq validiert\grqq{}. \emph{Viertens} regelt ein Fairness-Protokoll die
Praxis: Ein Verdikt, das vom eigenen Wortlaut des auditierten Autors abhängt,
führt diesen Wortlaut als wörtliches Zitat mit einer Adresse aus einer
benannten Version mit; jedes andere im Fließtext begründete Verdikt wird
durch eine Paraphrase der Passage oder des Merkmals getragen, auf das in der
benannten Version zurückgegriffen wird, oder durch eine Zitation in den
publizierten Serienteil, der es dokumentiert; die Zeilencode-Matrizen sind
Zusammenfassungen archivierter Auditnachweise, deren Zeilen-Ankertabellen
noch nicht publiziert sind -- bis das geplante Supplement erscheint, sind
die Matrizen als Zusammenfassungen jener Nachweise zu lesen und nicht als
unabhängig prüfbare Behauptungen (Abschnitt~\ref{sec:anchors}); die Konfidenz
wird in einem festen dreiachsigen Schema angegeben; Befunde der Form \glqq der
Text zeigt nicht\grqq{} werden niemals zu \glqq dies ist falsch\grqq{} aufgewertet; und die
eigenen Geltungsbereichsaussagen auditierter Autoren werden zitiert, nicht
paraphrasiert, wo immer ein Verdikt von ihnen abhängt.

Die zwei retrospektiven Fälle verwenden zudem unterschiedliche
Quellengattungen und auditierte Versionen: Der Wiles-Fall wird über den
zeitgenössischen Bericht von Rubin--Silverberg erschlossen, der Prasad-Fall
über das Primärpaper. Ihre rohen v1-T-Paare bleiben sichtbar: Das
Wiles-T-R8-Paar ist $(E-,I)$, das Prasad-T-R6-Paar $(I,E+)$.
Diese Quellenasymmetrie und die Übereinstimmung $10/12$ schätzen weder
Reliabilität noch Trennschärfe oder Fehlalarmrate; keines der Paare wird
stillschweigend in einen v1-F-F-R8-Code umgewandelt. Der LANA-Vergleich
verlangt weiterhin seine eigene Objekt-, Domänen-, Quantoren- und
Implikationsrichtungsbrücke; aus struktureller Ähnlichkeit folgt weder
Äquivalenz noch Bestätigung.

%=====================================================================
\section{Fazit}\label{sec:conclusion}

Die aktuelle Illustration erhält nun den vollständigen bedingten
C7-Vertrag; das unvollständige frühere Diagramm wurde korrigiert, ohne
Definition oder retrospektive Rohcodes zu ändern. Die aktuellen Teile 4--6
präzisieren die verbleibenden Normierungs- und Größenverbraucherpflichten,
ohne eine allgemeine Methodenvalidierung zu liefern oder Feld-, Perioden-,
Volumen- oder LANA-Transferbrücken zu schließen.



Wir haben ein prüfbares Formkriterium für Transportargumente angegeben:
sieben Vertragskomponenten, ein versioniertes Ledger mit einer
Drei-Verdikt-Grammatik über fünf Rohcodes und eine Auditpraxis mit wörtlicher
Verankerung und vor der Kodierung festgelegten Bewertungskriterien. Seine
Stellung unter bestehenden Strukturen ist spezifisch -- das nächste
heuristische Analogon ist die Tradition zertifizierender Algorithmen, aber
das Artefakt wird hier von einem dritten Auditor rekonstruiert statt vom
Beweisenden ausgestellt, und was eine Prüfung feststellt, ist textliche
Vertragsvollständigkeit und nicht Korrektheit --, und innerhalb des hier
gesichteten begrenzten Korpus fanden wir keine Arbeit, die in der Literatur
zu mathematischen Kontroversen denselben Platz einnimmt.

Der hier erbrachte Beitrag ist eine abc-Instanziierung eines
Form-Audit-Kriteriums. Die beiden historischen Fälle illustrieren v1-T
retrospektiv; sie validieren nicht das volle v1-F-Instrument. RH- und
global-only-Anwendungen bleiben künftige Arbeit. In der
protokoll-eingefrorenen Zwei-Fall-Illustration wurden die vorab festgelegten
Kriterien erfüllt. Angewandt auf den ausgewählten abc-Transportschritt
ersetzt das Ledger eine diffuse Uneinigkeit über diesen Schritt durch ein
Profil benannter Pflichten, gewichtet und jeweils mit einer Zitation; es
entscheidet nicht die Kontroverse, die ihn umgibt. Was das Kriterium nicht
kann, ist irgendetwas beweisen, und eine bediente Zeile zertifiziert keine
Mathematik. Sein Anspruch ist kleiner und, wie wir glauben, nützlich: dass
Uneinigkeiten über Transportbehauptungen als Uneinigkeiten über benannte
Vertragsklauseln geführt werden können -- prüfbar, zitierbar und
revidierbar, wenn sich der Text ändert.

%=====================================================================
% AI disclosure. The path is relative to the BUILD WORKING DIRECTORY, not to
% this file. Verified 2026-08-16 by compiling with cwd = paper/ , where it
% resolves to .RESEARCH/_templates/disclosure/ai_disclosure_STANDARD_de.tex and
% the disclosure is visible in the PDF. Compiling from any other directory
% requires adapting the path (or passing -include-directory).
% =============================================================================
% AI DISCLOSURE -- PIPELINE-STANDARD, DEUTSCHE FASSUNG (seit 2026-08-17, T-20260817-02)
% =============================================================================
% Sprachreine Fassung fuer rein deutschsprachige Manuskripte. Text 1:1 aus dem
% deutschen Block von ai_disclosure_STANDARD.tex uebernommen (keine inhaltliche
% Aenderung) -- nur die Ueberschrift ist jetzt sprachrein und der englische
% Block entfaellt. Die bilinguale Datei ai_disclosure_STANDARD.tex bleibt
% unveraendert bestehen (Bestandsverweise duerfen weiter darauf zeigen).
%
% ABWEICHEN NACH UNTEN: Wurde in einem Paper tatsaechlich deutlich weniger
% KI eingesetzt, wird NICHT auf eine andere Datei gewechselt, sondern die
% Formulierung hier abgeschwaecht -- z. B. "in erheblichem Masse in allen
% Phasen" ersetzen durch die tatsaechlich zutreffenden Phasen, und aus der
% Qualitaetssicherungs-Aufzaehlung streichen, was nicht stattfand.
% Der Grundsatz bleibt: ehrlich benennen, was war -- nicht abgrenzen.
%
% Verwendung (nur fuer deutschsprachige Manuskripte):
%   \input{../../_templates/disclosure/ai_disclosure_STANDARD_de}
% =============================================================================

\section*{KI-Offenlegung}

\noindent
KI-Systeme wurden in erheblichem Ma\ss{}e in allen Phasen dieser Arbeit
eingesetzt: Konzeption und Forschungsdiskussion, Literatur- und
Quellenarbeit, mathematische Analyse und Beweisarbeit, Erstellung und
Pr\"ufung der Rechenskripte, Ausf\"uhrung und Auswertung der
Berechnungsl\"aufe, Text, \"Ubersetzung und Satz.

\medskip

\noindent
Qualit\"atssicherung: Einsatz mehrerer unabh\"angiger Modelle
verschiedener Anbieter mit Gegenpr\"ufungen und adversarialen KI-Reviews;
nachvollziehbare, versionierte Rechenskripte mit gespeicherten
Ergebnisdateien und Pr\"ufsummen; fortlaufende Beweis- und
Registerdokumente; strikte Claim-Hygiene (Behauptungen nur mit
dokumentiertem Beleg, Nachtr\"age statt stiller Korrekturen); Changelogs
\"uber alle Versionen.


%=====================================================================
\begin{thebibliography}{99}
\sloppy
\emergencystretch=3em


\bibitem{Aberdein2023}
A.~Aberdein,
\emph{Deep Disagreement in Mathematics},
Global Philosophy \textbf{33} (2023), no.~1, article~17.
DOI: 10.1007/s10516-023-09653-7 (arXiv:2210.16488).

\bibitem{Andersen2017}
L.~E.~Andersen,
\emph{On the Nature and Role of Peer Review in Mathematics},
Accountability in Research \textbf{24} (2017), no.~3, 177--192.
DOI: 10.1080/08989621.2016.1274885.

\bibitem{Andersen2020}
L.~E.~Andersen,
\emph{Acceptable gaps in mathematical proofs},
Synthese \textbf{197} (2020), no.~1, 233--247.
DOI: 10.1007/s11229-018-1778-8.

\bibitem{GSN2021}
Assurance Case Working Group (ACWG),
\emph{Goal Structuring Notation Community Standard (Version 3)},
Safety-Critical Systems Club (SCSC), 4 May 2021.
Available at \url{https://scsc.uk/r1386.pdf}.

\bibitem{Avigad2021}
J.~Avigad,
\emph{Reliability of mathematical inference},
Synthese \textbf{198} (2021), no.~8, 7377--7399.
DOI: 10.1007/s11229-019-02524-y.

\bibitem{Bloomfield2021}
R.~Bloomfield and J.~Rushby, \emph{Assurance 2.0: A Manifesto},
arXiv:2004.10474v3 [cs.SE], 2020 (revised 14~January 2021).
DOI: 10.48550/arXiv.2004.10474.

\bibitem{Blum1995}
M.~Blum and S.~Kannan,
\emph{Designing programs that check their work},
Journal of the ACM \textbf{42} (1995), no.~1, 269--291.
DOI: 10.1145/200836.200880.

\bibitem{Buzzard2020}
K.~Buzzard, J.~Commelin and P.~Massot,
\emph{Formalising perfectoid spaces},
in: Proceedings of the 9th ACM SIGPLAN International Conference on
Certified Programs and Proofs (CPP 2020), ACM, 2020, 299--312.
DOI: 10.1145/3372885.3373830.

\bibitem{Castelvecchi2020}
D.~Castelvecchi,
\emph{Mathematical proof that rocked number theory will be published},
Nature \textbf{580} (2020), no.~7802, 177.
DOI: 10.1038/d41586-020-00998-2.

\bibitem{DeMillo1979}
R.~A.~De~Millo, R.~J.~Lipton and A.~J.~Perlis,
\emph{Social processes and proofs of theorems and programs},
Communications of the ACM \textbf{22} (1979), no.~5, 271--280.
DOI: 10.1145/359104.359106.

\bibitem{part1}
L.~Geiger,
\emph{The Gap in IUTchIII, Corollary 3.12: From Attempted Repair to
Structural Diagnosis},
Zenodo, 2026.
Concept DOI: 10.5281/zenodo.19960781
(latest version v6.3, DOI: 10.5281/zenodo.21899747).

\bibitem{part3}
L.~Geiger,
\emph{The Bridge Certificate: A Line-Level Audit Instrument for
Inter-Theoretic Transfer Arguments, with Four Applications to the abc
Literature (DRAFT)},
Zenodo, 2026.
Concept DOI: 10.5281/zenodo.21925803
(version 0.1, DOI: 10.5281/zenodo.21925804).

\bibitem{part4}
L.~Geiger,
\emph{The Bridge Certificate Applied to the Joshi ATS Series:
A Ledger Audit (DRAFT)},
Zenodo, 2026.
Concept DOI: 10.5281/zenodo.21951155
(version 0.2, DOI: 10.5281/zenodo.21951156).

\bibitem{part5}
L.~Geiger,
\emph{The Joshi ATS Series under the Bridge-Certificate Standard:
A Targeted Substance Audit and an All-Place Rigidity Classification of
Effective Normalization Weights (DRAFT)},
Zenodo, 2026.
Concept DOI: 10.5281/zenodo.21956398
(version 0.2, DOI: 10.5281/zenodo.21956399).

\bibitem{part6}
L.~Geiger,
\emph{Repairing the ATS Chain? A Rigidity Theorem for Product-Formula
Weights and an Effective-Weight Test for the Height Subchain (DRAFT)},
Zenodo, 2026.
Concept DOI: 10.5281/zenodo.21952486
(version 0.2, DOI: 10.5281/zenodo.21952487).

\bibitem{part4current}
L.~Geiger,
\emph{The Bridge Certificate Applied to the Joshi ATS Series: A Ledger Audit},
Zenodo, version 0.3, 2026. DOI: 10.5281/zenodo.23092802.

\bibitem{part5current}
L.~Geiger,
\emph{The Joshi ATS Series under the Bridge-Certificate Standard:
A Targeted Substance Audit and an All-Place Rigidity Classification
of Effective Normalization Weights},
Zenodo, version 0.3, 2026. DOI: 10.5281/zenodo.23094093.

\bibitem{part6current}
L.~Geiger,
\emph{What the Product Formula Allows: A Rigidity Theorem and
an Effective-Weight Test for the ATS Height Subchain},
Zenodo, version 0.3, 2026. DOI: 10.5281/zenodo.23093243.

\bibitem{gnc}
L.~Geiger,
\emph{Global Normalization in Joshi's Arithmetic Teichmuller Theory:
What the Identification Carries --- and What Carries the Theorem
(DRAFT)},
Zenodo, 2026.
Concept DOI: 10.5281/zenodo.21924253
(version 1.1, DOI: 10.5281/zenodo.21924254).

\bibitem{landscapeA}
L.~Geiger,
\emph{The abc Landscape: Obstructions, Failed Routes, and HCT
Diagnostics (DRAFT)},
Zenodo, 2026.
Concept DOI: 10.5281/zenodo.21916900
(latest version 1.1, DOI: 10.5281/zenodo.21924338).

\bibitem{zookeeper}
L.~Geiger,
\emph{The Spectral Zookeeper: A Conditional Reduction toward the
Riemann Hypothesis via Spectral Microcluster Closure in the
Connes--Consani--Moscovici Framework},
Zenodo, 2026.
Concept DOI: 10.5281/zenodo.19673126
(latest version 1.6, DOI: 10.5281/zenodo.21953305).

\bibitem{Hales2017}
T.~Hales et al.,
\emph{A formal proof of the Kepler conjecture},
Forum of Mathematics, Pi \textbf{5} (2017), e2.
DOI: 10.1017/fmp.2017.1.

\bibitem{JoshiATSIII}
K.~Joshi,
\emph{Construction of Arithmetic Teichmuller Spaces III: A `Rosetta
Stone' and a proof of Mochizuki's Corollary 3.12},
arXiv:2401.13508v4 [math.AG], 2024 (revised 24 February 2025).
DOI: 10.48550/arXiv.2401.13508.

\bibitem{JoshiATSIV}
K.~Joshi,
\emph{Construction of Arithmetic Teichmuller Spaces IV: Proof of the
abc-conjecture},
arXiv:2403.10430v2 [math.AG], 2024 (revised 24 February 2025).
The title page designates the text a \emph{preliminary version for
comments}.
DOI: 10.48550/arXiv.2403.10430.

\bibitem{Joshi2025Report}
K.~Joshi,
\emph{Final Report on the Mochizuki--Scholze--Stix Controversy},
arXiv:2505.10568v1 [math.AG], 2025 (submitted 29 April 2025), 8~pp.
DOI: 10.48550/arXiv.2505.10568.

\bibitem{LANA2026}
LANA Project,
\emph{Project LANA Interim Report on IUT Theory},
16 July 2026, 50~pp.
Available at
\url{https://ncatlab.org/nlab/files/LANAProject-Report-July2026.pdf}.

\bibitem{Mathlib2020}
The mathlib Community,
\emph{The Lean Mathematical Library},
in: Proceedings of the 9th ACM SIGPLAN International Conference on
Certified Programs and Proofs (CPP 2020), ACM, 2020, 367--381.
DOI: 10.1145/3372885.3373824.

\bibitem{McConnell2011}
R.~M.~McConnell, K.~Mehlhorn, S.~Näher and P.~Schweitzer,
\emph{Certifying algorithms},
Computer Science Review \textbf{5} (2011), no.~2, 119--161.
DOI: 10.1016/j.cosrev.2010.09.009.

\bibitem{Mochizuki2019}
S.~Mochizuki,
\emph{Report on Discussions, Held during the Period March 15--20, 2018,
Concerning Inter-universal Teichmüller Theory (IUTch)},
February 2019.
Available at
\url{https://www.kurims.kyoto-u.ac.jp/~motizuki/Rpt2018.pdf}.

\bibitem{Mochizuki2021c}
S.~Mochizuki,
\emph{Inter-universal Teichmüller Theory III: Canonical Splittings
of the Log-Theta-Lattice},
Publ.\ Res.\ Inst.\ Math.\ Sci.\ \textbf{57} (2021), no.~1/2, 403--626.
DOI: 10.4171/PRIMS/57-1-3.

\bibitem{Mochizuki2021d}
S.~Mochizuki,
\emph{Inter-universal Teichmüller Theory IV: Log-Volume
Computations and Set-Theoretic Foundations},
Publ.\ Res.\ Inst.\ Math.\ Sci.\ \textbf{57} (2021), no.~1/2, 627--723.
DOI: 10.4171/PRIMS/57-1-4.

\bibitem{Necula1997}
G.~C.~Necula,
\emph{Proof-carrying code},
in: Proceedings of the 24th ACM SIGPLAN-SIGACT Symposium on Principles
of Programming Languages (POPL '97), ACM, 1997, 106--119.
DOI: 10.1145/263699.263712.

\bibitem{prasad}
G.~Prasad,
\emph{Volumes of S-arithmetic quotients of semi-simple groups},
Publications Math\'{e}matiques de l'IH\'{E}S \textbf{69} (1989),
91--114.
DOI: 10.1007/BF02698841.
Available at \url{http://www.numdam.org/item/PMIHES_1989__69__91_0/}.

\bibitem{rubinsilverberg}
K.~Rubin and A.~Silverberg,
\emph{A report on Wiles' Cambridge lectures},
Bull.\ Amer.\ Math.\ Soc.\ (N.S.) \textbf{31} (1994), no.~1, 15--38.
DOI: 10.1090/S0273-0979-1994-00512-9 (arXiv:math/9407220).

\bibitem{Rushby2015}
J.~Rushby,
\emph{The Interpretation and Evaluation of Assurance Cases},
Technical Report SRI-CSL-15-01, Computer Science Laboratory,
SRI International, Menlo Park, CA, July 2015.
Available at
\url{https://www.csl.sri.com/users/rushby/papers/sri-csl-15-1-assurance-cases.pdf}.

\bibitem{Scholze2022LTE}
P.~Scholze,
\emph{Liquid Tensor Experiment},
Experimental Mathematics \textbf{31} (2022), no.~2, 349--354.
DOI: 10.1080/10586458.2021.1926016.

\bibitem{ScholzeStix2018}
P.~Scholze and J.~Stix,
\emph{Why abc is still a conjecture},
unpublished manuscript, 2018; the circulated file was produced on
16 July 2018.
Available at
\url{https://www.math.uni-bonn.de/people/scholze/WhyABCisStillaConjecture.pdf}.

\bibitem{ZhouI}
Z.-P.~Zhou,
\emph{The inter-universal Teichmüller theory and new Diophantine
results over the rational numbers. I},
arXiv:2503.14510v1 [math.NT], 2025.
DOI: 10.48550/arXiv.2503.14510.

\bibitem{ZhouII}
Z.-P.~Zhou,
\emph{The inter-universal Teichmüller theory and new Diophantine
results over rational numbers. II},
arXiv:2510.05448v1 [math.NT], 2025.
The posting designates itself an older version awaiting update.
DOI: 10.48550/arXiv.2510.05448.

\end{thebibliography}

\end{document}
